% Modern editorial, markup and figure/layout contributions: CC-BY-SA-4.0.
% Historical text and illustrations: Leonhard Euler (1744), public domain.
% Component scope and upstream attribution: see RIGHTS.md.
\subsection*{Additamentum I.}
\begin{center}\textit{De curvis elasticis.}\end{center}
% Additamentum I. Image-collated transcription. PDF 248--313.
\originalsection{1}\sourcepage{245}Jam pridem summi quique Geometræ agnoverunt, Methodi in hoc Libro traditæ non solum maximum esse usum in ipsa Analysi, sed etiam eam ad resolutionem Problematum physicorum amplissimum subsidium afferre. Cum enim Mundi universi fabrica sit perfectissima, atque a Creatore sapientissimo absoluta, nihil omnino in mundo contingit, in quo non maximi minimive ratio quæpiam eluceat: quamobrem dubium prorsus est nullum, quin omnes Mundi effectus ex causis finalibus, ope Methodi maximorum \& minimorum æque feliciter determinari queant, atque ex ipsis causis efficientibus. Hujus rei vero passim tam eximia extant specimina, ut ad veritatis confirmationem pluribus Exemplis omnino non indigeamus; quin potius in hoc erit elaborandum, ut, in quovis Quæstionum naturalium genere, ea investigetur quantitas, quæ maximum minimumve induat valorem: quod negotium ad Philosophiam potius quam ad Mathesin pertinere videtur. Cum igitur duplex pateat via effectus Naturæ cognoscendi; altera per causas efficientes, quæ Methodus directa vocari solet; altera causas finales; Mathematicus utrâque pari successu utitur. Quando scilicet causæ efficientes nimis sunt absconditæ, finales autem nostram cognitionem minus effugiunt; per Methodum indirectam Quæstio solet resolvi: e contrario autem Methodus directa adhibetur, quoties ex causis efficientibus effectum definire licet. Inprimis autem opera est adhibenda, ut per utramque viam aditus ad Solutionem aperiatur: sic enim non solum altera Solutio per alteram maxime confirmatur, sed etiam ex utriusque consensu \sourcepage{246}summam percipimus voluptatem. Hoc modo, curvatura funis seu catenæ suspensæ duplici via est eruta; altera a priori, ex sollicitationibus gravitatis; altera vero per Methodum maximorum ac minimorum, quoniam funis ejusmodi curvaturam recipere debere intelligebatur, cujus centrum gravitatis infimum obtineret locum. Similiter curvatura radiorum per medium diaphanum variæ densitatis transeuntium, tam a priori est determinata, quam etiam ex hoc principio, quod tempore brevissimo ad datum locum pervenire debeant. Plurima autem alia similia exempla a Viris Celeberrimis Bernoulliis, aliisque, sunt prolata, quibus tam Methodus solvendi a priori, quam cognitio causarum efficientium maxima accepit incrementa. Quanquam igitur, ob hæc tam multa ac præclara specimina, dubium nullum relinquitur, quin in omnibus lineis curvis, quas Solutio Problematum physico-mathematicorum suppeditat, maximi minimive cujuspiam indoles locum obtineat; tamen sæpenumero hoc ipsum maximum vel minimum difficillime perspicitur; etiamsi a priori Solutionem eruere licuisset. Sic etsi figura, quam lamina elastica incurvata induit, jam pridem est cognita; tamen quemadmodum ea curva per Methodum maximorum \& minimorum, hoc est, per causas finales, investigari possit, a nemine adhuc est animadversum. Quamobrem cum Vir Celeberrimus, atque in hoc sublimi naturam scrutandi genere perspicacissimus, \textit{Daniel Bernoulli} mihi indicasset se universam vim, quæ in lamina elastica incurvata insit, una quadam formula quam \textit{vim potentialem} appellat, complecti posse; hancque expressionem in curva Elastica minimam esse oportere; quoniam hoc invento Methodus mea maximorum ac minimorum hoc Libro tradita mirifice illustratur, ejusque usus amplissimus maxime evincitur; hanc occasionem exoptatissimam prætermittere non possum, quin, hanc insignem curvæ Elasticæ proprietatem a Celeb. Bernoullio observatam publicando, simul Methodi meæ usum clarius patefaciam. Continet enim ista proprietas in se differentialia secundi gradus, ita ut ei evolvendæ Methodi Problema isoperimetricum solvendi ante traditæ non sufficiant.

\originalsection{2}\sourcepage{247}\marginnote{\hyperlink{addfig:1}{Fig.~1}}[1.2\baselineskip]Sit $AB$ lamina Elastica utcunque incurvata; vocetur arcus $AM=s$, \& radius osculi curvæ $MR=R$: atque, secundum Bernoullium, exprimetur \textit{vis potentialis} in laminæ portione $AM$ contenta hac formula $\int\frac{ds}{RR}$, siquidem lamina sit ubique æqualiter crassa, lata \& elastica, atque in statu naturali in directum extensa. Hinc ista erit curvæ $AM$ indoles, ut in ea hæc expressio omnium minimum obtineat valorem. Quoniam vero in radio osculi $R$ differentialia secundi gradus insunt, ad curvam hac proprietate præditam determinandam quatuor opus erit conditionibus, id quod cum Quæstionis natura apprime convenit. Cum enim per datos terminos $A$ \& $B$ infinitæ laminæ Elasticæ eæque ejusdem longitudinis inflecti queant, quæstio non erit determinata, nisi præter duo puncta $A$ \& $B$, simul alia duo puncta, seu quod eodem redit positio tangentium in punctis extremis $A$ \& $B$ præscribatur. Proposita namque lamina Elastica, longiori quam est distantia punctorum $A$ \& $B$; ea non solum ita incurvari potest, ut intra terminos $A$ \& $B$ contineatur, sed etiam ut ejus tangentes in punctis hisce datas teneant directiones. His notatis; Quæstio de invenienda curvatura laminæ Elasticæ, ex hoc fonte resolvenda, ita debet proponi: \textit{ut, inter omnes curvas ejusdem longitudinis, quæ non solum per puncta $A$ \& $B$ transeant, sed etiam in his punctis a rectis positione datis tangantur, definiatur ea in qua sit valor hujus expressionis $\int\frac{ds}{RR}$ minimus.}

\originalsection{3}\marginnote{\hyperlink{addfig:2}{Fig.~2}}Quia solutionem ad coordinatas orthogonales accommodari convenit, sumatur recta quæcunque $AD$ pro axe, in qua sit abscissa $AP=x$, applicata $PM=y$, ponatur, uti Methodus tradita jubet, $dy=pdx$, $dp=qdx$; erit elementum curvæ $Mm=ds=dx\sqrt{1+pp}$. Primum ergo quia curvæ, ex quibus quæsita erui debet, isoperimetræ statuuntur, habebitur ista expressio consideranda $\int dx\sqrt{1+pp}$; quæ cum generali $\int Zdx$ comparata hunc præbet valorem differentialem
\sourcepage{248}
\[\frac{1}{dx}d.\frac{p}{\sqrt{1+pp}}.\]
Deinde cum sit radius osculi $\frac{dx(1+pp)^{3:2}}{dp}=\frac{(1+pp)^{3:2}}q=R$, formula $\int\frac{ds}{RR}$, quæ minimum esse debet, abit in $\int\frac{qqdx}{(1+pp)^{5:2}}$. Comparetur hæc cum forma generali $\int Zdx$; erit $Z=\frac{qq}{(1+pp)^{5:2}}$, \& posito $dZ=Mdx+Ndy+Pdp+Qdq$, erit
\[M=0,\quad N=0,\quad P=\frac{-5pqq}{(1+pp)^{7:2}},\quad\&\quad Q=\frac{2q}{(1+pp)^{5:2}}.\]
Valor ergo differentialis ex hac formula $\int\frac{qqdx}{(1+pp)^{5:2}}$ oriundus, erit $-\frac{dP}{dx}+\frac{ddQ}{dx^2}$. Quamobrem pro curva quæsita hæc habebitur æquatio,
\[\frac{\alpha}{dx}d.\frac p{\sqrt{1+pp}}=\frac{dP}{dx}-\frac{ddQ}{dx^2};\]
quæ, per $dx$ multiplicata \& integrata, dat
\[\frac{\alpha p}{\sqrt{1+pp}}+\beta=P-\frac{dQ}{dx}.\]
Multiplicetur hæc æquatio per $qdx=dp$, ut prodeat
\[\frac{\alpha p\,dp}{\sqrt{1+pp}}+\beta dp=Pdp-qdQ.\]
Cum autem, ob $M=0$ \& $N=0$, sit $dZ=Pdp+Qdq$, erit $Pdp=dZ-Qdq$, quo valore loco $Pdp$ substituto, emerget
\[\frac{\alpha p\,dp}{\sqrt{1+pp}}+\beta dp=dZ-Qdq-qdQ;\]
quæ denuo integrata dat $\alpha\sqrt{1+pp}+\beta p+\gamma=Z-Qq$. Jam cum sit $Z=\frac{qq}{(1+pp)^{5:2}}$, \& $Q=\frac{2q}{(1+pp)^{5:2}}$, erit
\[\alpha\sqrt{1+pp}+\beta p+\gamma=\frac{-qq}{(1+pp)^{5:2}}.\]
Sumantur constantes arbitrariæ $\alpha$, $\beta$, \& $\gamma$ negative, eritque
\sourcepage{249}
\[q=(1+pp)^{5:4}\sqrt{\alpha\sqrt{1+pp}+\beta p+\gamma}=\frac{dp}{dx}.\]
Hinc ergo elicitur sequens æquatio
\[dx=\frac{dp}{(1+pp)^{5:4}\sqrt{\alpha\sqrt{1+pp}+\beta p+\gamma}}.\]
Deinde ob $dy=pdx$, habebitur quoque
\[dy=\frac{p\,dp}{(1+pp)^{5:4}\sqrt{\alpha\sqrt{1+pp}+\beta p+\gamma}};\]
quæ duæ æquationes sufficient ad curvam per quadraturas construendam.

\originalsection{4}Harum formularum sic in genere spectatarum neutra est integrabilis; combinari autem certo quodam modo possunt, ut aggregatum integrationem admittat. Cum enim sit
\[d.\frac{2\sqrt{\alpha\sqrt{1+pp}+\beta p+\gamma}}{\sqrt{\sqrt{1+pp}}}=\frac{dp(\beta-\gamma p)}{(1+pp)^{5:4}\sqrt{\alpha\sqrt{1+pp}+\beta p+\gamma}},\]
erit $\frac{2\sqrt{\alpha\sqrt{1+pp}+\beta p+\gamma}}{(1+pp)^{1:4}}=\beta x-\gamma y+\delta$. Quoniam axis positio est arbitraria, constans $\delta$ sine defectu amplitudinis omitti potest. Deinde vero etiam axis ita mutari potest ut fiat $\frac{\beta x-\gamma y}{\sqrt{\beta\beta+\gamma\gamma}}$ abscissa, eritque applicata $\frac{\gamma x+\beta y}{\sqrt{\beta\beta+\gamma\gamma}}$; hinc etiam tuto $\gamma$ nihilo æqualis poni potest, quia nihil impedit, quominus illa nova abscissa per $x$ exprimatur. Hanc ob rem; habebimus pro curva Elastica istam æquationem
\[2\sqrt{\alpha\sqrt{1+pp}+\beta p}=\beta x(1+pp)^{1:4};\]
quæ, sumptis quadratis, dat $4\alpha\sqrt{1+pp}+4\beta p=\beta^2x^2\sqrt{1+pp}$. Sit, ad homogeneitatem introducendam, $\alpha=\frac{4m}{aa}$ \& $\beta=\frac{4n}{aa}$ erit $naap=(nnxx-maa)\sqrt{1+pp}$, unde $n^2a^4pp=(nnxx-maa)^2(1+pp)$; ideoque
\[p=\frac{nnxx-maa}{\sqrt{n^2a^4-(nnxx-maa)^2}}=\frac{dy}{dx}.\]
Mutatis ergo constantibus, atque abscissam $x$ data constante sive augendo sive minuendo; habebitur hujusmodi æquatio pro curva Elastica generalis:
\sourcepage{250}
\[dy=\frac{(\alpha+\beta x+\gamma xx)dx}{\sqrt{a^4-(\alpha+\beta x+\gamma xx)^2}},\]
ex qua oritur
\[ds=\frac{aa\,dx}{\sqrt{a^4-(\alpha+\beta x+\gamma xx)^2}};\]
ex quibus æquationibus consensus hujus curvæ inventæ cum curva Elastica jam pridem eruta manifesto elucet.

\originalsection{5}Quo autem iste consensus clarius ob oculos ponatur, naturam curvæ Elasticæ a priori quoque investigabo; quod etsi jam a Viro summo \textit{Jacobo Bernoullio} excellentissime est factum; tamen, hac idonea occasione oblata, nonnulla circa indolem curvarum Elasticarum, earumque varias species \& figuras adjiciam; quæ ab aliis vel prætermissa, vel leviter tantum pertractata esse video.

\marginnote{\hyperlink{addfig:3}{Fig.~3}}Sit lamina Elastica $AB$ in $B$ ita muro seu pavimento firmo infixa, ut hæc extremitas $B$ non solum firmiter retineatur, sed etiam tangentis in $B$ positio determinetur. In $A$ autem lamina connexam habeat virgam rigidam $AC$, cui normaliter applicata sit vis $CD=P$, qua lamina in statum incurvatum $BMA$ redigatur. Sumatur hæc recta $AC$ producta pro axe, ac, posita $AC=c$, sit abscissa $AP=x$, applicata $PM=y$. Quod si jam lamina in $M$ omnem elasticitatem subito amitteret, ac perfecte flexilis evaderet; a vi $P$ utique inflecteretur, inflexione proficiscente a vis $P$ momento $=P(c+x)$. Quominus ergo hæc inflexio actu sequatur, elasticitas laminæ in $M$ in æquilibrio consistit cum vis sollicitantis momento $P(c+x)$. Elasticitas autem primo ab indole materiæ ex qua lamina constat, \& quam ubique eandem statuo, pendet; tum vero simul ab incurvatione laminæ in puncto $M$, ita ut sit reciproce proportionalis radio osculi in $M$. Sit ergo radius osculi in $M=R=\frac{ds^3}{-dx\,ddy}$; existente $ds=\sqrt{dx^2+dy^2}$ \& $dx$ constante; atque exprimat $\frac{Ekk}{R}$ vim Elasticam laminæ in $M$, quæ cum momento vis sollicitantis $P(c+x)$ in æquilibrio consistat, ita ut sit
\[P(c+x)=\frac{Ekk}{R}=\frac{-Ekk\,dx\,ddy}{ds^3}.\]
Æquatio hæc
\sourcepage{251}per $dx$ multiplicata fit integrabilis, eritque integrale
\[P(xx+cx+f)=\frac{-Ekk\,dy}{\sqrt{dx^2+dy^2}};\]
unde oritur
\[dy=\frac{-Pdx(\frac12xx+cx+f)}{\sqrt{E^2k^4-P^2(\frac12xx+cx+f)^2}},\]
quæ æquatio omnino convenit cum ea, quam modo per Methodum maximorum ac minimorum ex principio \textit{Bernoulliano} elicui.

\originalsection{6}Ex comparatione hujus æquationis cum ante inventa, definiri poterit vis quæ requiritur ad datam laminæ curvaturam inducendam; siquidem curvatura contineatur in æquatione generali inventa. Teneat scilicet lamina elastica figuram $AMB$, cujus natura exprimatur hac æquatione
\[dy=\frac{(\alpha+\beta x+\gamma xx)dx}{\sqrt{a^4-(\alpha+\beta x+\gamma xx)^2}};\]
exprimat vero $Ekk$ hujus laminæ elasticitatem absolutam, ita scilicet, ut $Ekk$, in quovis loco, per radium osculi divisa præbeat vim elasticam veram. Ad comparationem instituendam multiplicetur numerator \& denominator per $\frac{Ekk}{aa}$, ut habeatur
\[dy=\frac{Ekk\,dx(\alpha+\beta x+\gamma xx):aa}{\sqrt{E^2k^4-\frac{E^2k^4}{a^4}(\alpha+\beta x+\gamma xx)^2}}.\]
Nunc ergo erit $-\frac12P=\frac{Ekk\gamma}{aa}$; $-Pc=\frac{Ekk\beta}{aa}$; $-Pf=\frac{Ekk\alpha}{aa}$; ideoque vis $CD$ sollicitans $=\frac{-2Ekk\gamma}{aa}$; intervallum $AC=c=\frac\beta{2\gamma}$, \& constans $f=\frac\alpha{2\gamma}$.

\originalsection{7}Ut igitur lamina elastica $AB$ altero termino $B$ muro infixa incurvetur in figuram $AMB$, cujus natura exprimitur hac æquatione $dy=\frac{(\alpha+\beta x+\gamma xx)dx}{\sqrt{a^4-(\alpha+\beta x+\gamma xx)^2}}$, necesse est ut hæc lamina sollicitetur in directione $CD$ normali ad axem $AP$, sumpta distantia $AC=\frac\beta{2\gamma}$, a vi $CD=\frac{-2Ekk\gamma}{aa}$; quæ vis scilicet in plagam contrariam, ac figura indicat, dirigetur, \sourcepage{252}si $\gamma$ fuerit quantitas positiva. Quia $\frac{Ekk}{R}$ æquivalet momento vis sollicitantis, expressio $\frac{Ekk}{aa}$ homogenea erit ponderi seu vi puræ; quæ vis propterea $\frac{Ekk}{aa}$ cognoscetur ex elasticitate laminæ. Sit hæc vis $=F$; atque erit vis flectens $CD$ ad hanc vim $F$ ut $-2\gamma$ ad $1$; erit enim $\gamma$ numerus purus.

\originalsection{8}Hinc porro definiri potest vis ad laminæ portionem $BM$ in statu suo conservandam requisita, si portio $AM$ prorsus rescindatur. Rescissa hac portione $AM$, desinat lamina Elastica in virgam rigidam $MT$ omnis flexionis expertem, quæ autem cum lamina ita sit connexa, ut perpetuo tangentem in puncto $M$ referat, utcunque lamina inclinetur. Hoc posito, ex antecedentibus manifestum est, ad conservationem curvaturæ $BM$ requiri ut virga $MT$ in puncto $N$ trahatur in directione $ND$ vi quæ sit $=\frac{-2Ekk\gamma}{aa}$; directio autem $ND$ erit normalis ad axem $AP$, atque intervallum $AC$ erit $=\frac\beta{2\gamma}$. Distantia itaque $MN$ fiet $=\frac{ds}{dx}CP=\frac{ds}{dx}.\frac{\beta+2\gamma x}{2\gamma}=\frac{(\beta+2\gamma x)ds}{2\gamma dx}$ est vero $\frac{ds}{dx}=\frac{aa}{\sqrt{a^4-(\alpha+\beta x+\gamma xx)^2}}$. Quod si hæc vis $ND=\frac{-2Ekk\gamma}{aa}$ resolvatur in normalem $NQ$ ad tangentem $MT$, \& tangentialem $NT$, erit vis normalis $NQ=\frac{-2Ekk\gamma}{aa}.\frac{dx}{ds}$, \& vis tangentialis $NT=\frac{-2Ekk\gamma}{aa}.\frac{dy}{ds}$.

\originalsection{9}Sin autem pars $BM$ rescindatur, relicta parte $AM$, quæ in directione $CD$ sollicitatur ut ante vi $=\frac{-2Ekk\gamma}{aa}$; ad curvaturam $AM$ conservandam extremitas $M$, quæ connexa intelligatur cum virga rigida tangente $MN$, sollicitari debet in puncto $N$ a vi pariter $=\frac{-2Ekk\gamma}{aa}$, sed in directione contraria
\sourcepage{253}ei, quam casu præcedente invenimus. Perpetuo enim vires utrique extremitati laminæ incurvatæ applicandæ se mutuo destruere, atque adeo æquales \& directiones oppositas habere debent. Alioquin enim tota lamina moveretur, ad quem motum compescendum opus foret vi æquilibrium inter vires sollicitantes producente. Hinc ergo vires cuicunque portioni laminæ resectæ applicandæ facillime definiri possunt, quæ jam inductam curvaturam conservent.

\originalsection{10}\marginnote{\hyperlink{addfig:4}{Fig.~4}}Sit $AM$ lamina Elastica incurvata, quæ in $A$ \& $M$ annexas habeat virgas rigidas $AD$, $MN$, quibus in directionibus directe oppositis $DE$, $NR$ applicatæ sint vires æquales $DE$, $NR$, quæ in æquilibrio consistentes laminæ curvaturam $AM$ inducant; pro qua æquationem quæri oporteat. Primum ergo, pro axe sumatur recta $AP$ per punctum $A$ transiens, atque ad directionem vis sollicitantis $ER$ normalis. Ponatur Elasticitas laminæ absoluta $=Ekk$: sitque anguli $CAD$, quem tangens $AD$ in $A$ cum axe constituit, \& qui est datus, sinus $=m$, cosinus $=n$, existente sinu toto $=1$, ita ut sit $mm+nn=1$. Vocetur porro distantia $AC=c$, \& vis flectens $DE=NR=P$; ac, positis abscissa $AP=x$, applicata $PM=y$, natura curvæ hac exprimetur æquatione
\[dy=\frac{-Pdx(\frac12xx+cx+f)}{\sqrt{E^2k^4-P^2(\frac12xx+cx+f)^2}}.\]
Quoniam vero directio tangentis in $A$ datur, posito $x=0$, fieri debet $\frac{dy}{dx}=\frac mn$; hinc ergo obtinebitur $\frac mn=\frac{-Pf}{\sqrt{E^2k^4-P^2f}}=\frac m{\sqrt{1-mm}}$, \& $m=\frac{-Pf}{Ekk}$. Determinatur ergo hinc constans $f$, ita ut sit $f=\frac{-mEkk}{P}$, ideoque hinc tota curva determinatur.

\originalsection{11}\marginnote{\hyperlink{addfig:5}{Fig.~5}}Ad curvaturam ergo superiori æquatione expressam laminæ $AM$ inducendam, tangenti $AD$ in puncto $D$, ita ut sit $AD=\frac cn$, applicatam esse oportet vim $DE=P$; cujus \sourcepage{254}directio sit parallela applicatis $PM$. Resolvatur hæc vis $DE$ in duas laterales $Dd$, $Df$, inter se normales; erit vis $Dd=Pn$ \& vis $Df=Pm$. Quo jam consideratio rectæ $AD$ ex computo expellatur, loco vis $Dd$, in datis punctis $A$ \& $B$, sumpto intervallo $AB=h$, duæ vires substitui possunt $Aa=p$; $Bb=q$; normales pariter ad virgam $AB$, sumendo $ph=Pn.BD=nP(\frac cn-h)$, \& $q=p+nP$. Quia deinceps perinde est, in quonam virgæ $AD$ puncto applicetur vis tangentialis $Df=mP$, applicetur ea in ipso puncto $A$ ponendo $AF=nP$. Sit autem hæc vis $AF=r$, ita ut lamina $MA$ a tribus viribus $Aa=p$, $Bb=q$, \& $AF=r$ sollicitetur, a quibus, qualis incurvatio oriatur, investigemus.

\originalsection{12}Primo ergo, cum sit $mP=r$, erit $P=\frac rm$, qui valor substitutus in prioribus æquationibus dabit $ph=\frac{cr}m-\frac{nhr}m$, \& $q=p+\frac{nr}m$. Hinc erit $\frac nm=\frac{q-p}r$; ex qua æquatione primum positio axis $AP$ innotescit; erit nempe tangens anguli $CAD=\frac r{q-p}$: hinc $m=\frac r{\sqrt{r^2+(q-p)^2}}$, \& $n=\frac{q-p}{\sqrt{r^2+(q-p)^2}}$. Deinde ex æquatione $hp=\frac{cr}m-\frac{nhr}m=\frac{cr}m-hq+hp$; fit $c=\frac{mhq}r$, seu $c=\frac{hq}{\sqrt{r^2+(q-p)^2}}$; atque $P=\sqrt{rr+(q-p)^2}$. Cum autem sit $f=\frac{-mEkk}P=\frac{-Ekkr}{rr+(q-p)^2}$ erit
\[\frac12xx+cx+f=\frac12xx+\frac{hqx}{\sqrt{r^2+(q-p)^2}}-\frac{Ekkr}{rr+(q-p)^2};\]
unde pro curva quæsita ista obtinebitur æquatio
\[dy=\frac{dx\left(\frac{Ekkr}{\sqrt{rr+(q-p)^2}}-hqx-\frac12xx\sqrt{rr+(q-p)^2}\right)}{\sqrt{E^2k^4-\left(\frac{Ekkr}{\sqrt{rr+(q-p)^2}}-hqx-\frac12xx\sqrt{rr+(q-p)^2}\right)^2}}.\]
Hæc autem æquatio maxime est accommodata ad modum maxime
\sourcepage{255}consuetum laminas incurvandi, dum eæ vel forcipe vel duobus digitis apprehenduntur; quorum alter laminam in directione $Aa$, alter in directione $Bb$ urget, præter quas vires lamina insuper in directione $AF$ protrahi potest.

\originalsection{13}Si vis tangentialis $AF=r$ evanescat; incidet axis $AP$ in ipsam tangentem $AF$, productam, eritque tum
\[dy=\frac{-dx(hqx+\frac12(q-p)xx)}{\sqrt{E^2k^4-(hqx+\frac12(q-p)xx)^2}}.\]
Sin autem vires normales $p$ \& $q$ fiant inter se æquales; erit axis $AP$ normalis ad tangentem $AF$, ob $n=0$; \& pro curva orietur hæc æquatio
\[dy=\frac{dx(Ekk-hqx-\frac12rxx)}{\sqrt{2Ekk(hqx+\frac12rxx)-(hqx-\frac12rxx)^2}}.\]
Hic si præterea fuerit $r=0$, ita ut lamina in punctis $A$ \& $B$ urgeatur a viribus æqualibus $Aa$, $Bb$, contrariis tantum, natura curvæ exprimetur hac æquatione $dy=\frac{dx(Ekk-hqx)}{\sqrt{hq(2Ekkx-hqxx)}}$, quæ integrata dat $y=\sqrt{\frac{2Ekkx-hqxx}{hq}}$; quæ est pro Circulo, lamina ergo hoc casu in arcum Circuli incurvatur, cujus radius erit $=\frac{Ekk}{hq}$.

\originalsection{14}Cum igitur videamus non solum Circulum in curvarum Elasticarum classe contineri, sed etiam in ipsis infinitam varietatem locum habere; operæ pretium erit hic enumerationem omnium variarum specierum in hoc curvarum genere contentarum instituere. Hoc enim modo non solum indoles harum curvarum penitius perspicietur; sed etiam, casu quocunque oblato, ex sola figura dijudicare licebit, ad quamnam speciem curva formata referri debeat. Eodem autem modo hic specierum diversitatem constituemus, quo vulgo linearum algebraicarum species, in dato ordine contentæ enumerari solent.

\originalsection{15}Æquatio generalis pro curvis Elasticis $dy=\frac{(\alpha+\beta x+\gamma xx)dx}{\sqrt{a^4-(\alpha+\beta x+\gamma xx)^2}}$, initio abscissarum in axe \sourcepage{256}per intervallum $\frac\beta{2\gamma}$ promoto, \& pro $\frac{aa}\gamma$ scribendo $aa$, seu ponendo $\gamma=1$, accipiet hanc formam simpliciorem:
\[dy=\frac{(\alpha+xx)dx}{\sqrt{a^4-(\alpha+xx)^2}}.\]
Quia vero est $a^4-(\alpha+xx)^2=(aa-\alpha-xx)(aa+\alpha+xx)$; ponatur $aa-\alpha=cc$, ut sit $\alpha=aa-cc$, atque æquatio transibit in hanc formam
\[dy=\frac{(aa-cc+xx)dx}{\sqrt{(cc-xx)(2aa-cc+xx)}}.\]
\marginnote{\hyperlink{addfig:6}{Fig.~6}}Qua æquatione exprimatur natura curvæ $AMC$, posita abscissa $AP=x$, \& applicata $PM=y$. Cum ergo sit $\beta=0$, directio vis laminam Elasticam incurvantis erit ad axem $AP$ in ipso puncto $A$ normalis, ideoque $AD$ repræsentabit directionem vis sollicitantis, quæ vis ipsa erit $=\frac{2Ekk}{aa}$, exprimente $Ekk$ elasticitatem absolutam.

\originalsection{16}Si ponatur $x=0$, erit $\frac{dy}{dx}=\frac{aa-cc}{c\sqrt{2aa-cc}}$; quæ expressio præbet tangentem anguli quem curva $AM$ in $A$ cum axe $AP$ constituit; cujus anguli sinus erit $=\frac{aa-cc}{aa}$. Quare si fuerit $aa=\infty$, lamina in puncto $A$ erit normalis ad axem $AP$, nullamque habebit curvaturam, propterea quod vis incurvans $\frac{2Ekk}{aa}$ evanescit. Casu ergo quo $a=\infty$, prodit laminæ figura naturalis, hoc est linea recta: quæ ergo primam speciem linearum Elasticarum constituit, quam repræsentabit recta $AB$ utrinque in infinitum producta.

\originalsection{17}Antequam reliquas species enumeremus, conveniet in genere circa figuram Elasticæ quasdam observationes instituere. Intelligitur autem angulus $PAM$, quem curva in $A$ cum axe $AP$ constituit, decrescere, quo minor evadat quantitas $aa$, hoc est quo magis vis incurvans $\frac{2Ekk}{aa}$ intendatur. Atque si evadat $aa=cc$, tum axis $AP$ ipse curvam in $A$ tanget. Quod si autem fuerit $aa<cc$, tum curva $AM$, quæ adhuc deorsum excurrebat, nunc sursum verget, quoad fiat $aa=\frac12cc$; quo casu tangens
\sourcepage{257}curvæ in rectam $Ab$ incidet. At si fiat $aa<\frac12cc$, tum angulus $PAM$ prorsus fiet imaginarius, ideoque in $A$ nulla existet curvæ portio, qui diversi casus specierum varietatem constituent.

\originalsection{18}Ex æquatione porro intelligitur, quia formam suam non mutat, si coordinatæ $x$ \& $y$ ambæ negativæ statuantur, curvam circa $A$ ramos habere similes \& æquales $AMC$ \& $Amc$ alternatim dispositos; ita ut in $A$ sit punctum flexus contrarii; unde cognita curvæ portione $AMC$, simul ejus continuatio $Amc$ ultra $A$ cognoscetur, quippe quæ illi est similis \& æqualis. Sic sumpta $Ap=AP$, erit quoque $pm=PM$. Recedendo autem ab $A$, curva utrinque magis ab axe reclinatur, donec sumpta abscissa $=AE=c$, applicata $EC$ curvam tangat; namque posito $x=c$, fit $\frac{dy}{dx}=\infty$. Perspicuum autem est abscissam $x$ ultra $AE=c$ excrescere non posse; alioquin enim fieret $\frac{dy}{dx}$ imaginarium; hinc ergo tota curva continebitur inter applicatas extremas $EC$ \& $ec$, ultra quos cancellos egredi non queat. Jam ergo generatim cognitos habemus binos curvæ ramos $AC$ \& $Ac$ utrinque ab $A$ usque ad cancellos protensos.

\originalsection{19}Videamus ergo quonam cursu curva ultra $C$ \& $c$ progrediatur. Hunc in finem sumamus rectam $CD$ ipsi $AE$ parallelam pro axe, ac ponamus has novas coordinatas $CQ=t$, $QM=u$; eritque $t+x=AE=CD=c$; \& $y+u=CE=AD=b$; unde fit $x=c-t$ \& $y=b-u$, seu $dy=-du$. His valoribus substitutis, orietur æquatio pro curva inter coordinatas $CQ=t$ \& $QM=u$, quæ erit
\[du=\frac{(aa-2ct+tt)dt}{\sqrt{t(2c-t)(2aa-2ct+tt)}}.\]
Hic primum patet, si sumatur $t$ infinite parvum, fore $du=\frac{aa\,dt}{2a\sqrt{ct}}$, ideoque $u=a\sqrt{\frac tc}$; quæ æquatio indicat curvam ultra $C$ simili modo versus \sourcepage{258}$N$ progredi incipere, quo ex $C$ ad $M$ extenditur. Ambiguitas autem signi $\sqrt{\vphantom a}$ in denominatore æquationis luculenter declarat, applicatam $u$ æque negative accipi posse atque affirmative: unde manifestum est, rectam $CD$ esse curvæ diametrum, atque adeo arcum $CNB$ similem \& æqualem fore arcui $CMA$.

\originalsection{20}Simili autem modo recta $cd$, ex altera parte axi $AE$ per $c$ parallela ducta, erit curvæ diameter; propterea quod ramus $Acb$ similis \& æqualis est ramo $ACB$. In punctis ergo $B$ \& $b$, erunt quoque puncta flexus contrarii omnino uti in $A$; unde curva similiter ulterius progredietur. Habebit ergo curva infinitas diametros $CD$, $cd$, \&c. intervallo eodem $Dd$ a se invicem distantes ac parallelas inter se; hancque ob rem curva constabit ex infinitis partibus inter se similibus \& æqualibus; atque ideo tota curva cognoscetur, si unica tantum portio $AMC$ fuerit perspecta.

\originalsection{21}Quia in $A$ est punctum flexus contrarii, ibidem erit radius osculi infinite magnus; id quod ex ipsa curvæ natura patet. Cum enim curva in $A$ sollicitetur a vi $=\frac{2Ekk}{aa}$ in directione $AD$; erit in quovis loco $M$, si radius osculi ibi ponatur $=R$, ex natura elasticitatis $\frac{2Ekk}{aa}x=\frac{Ekk}R$; unde fit $R=\frac{aa}{2x}$. In puncto ergo $A$ radius osculi est infinitus; at vero in punctis $C$, $c$, ob $AE=Ae=c$, erit radius osculi $=\frac{aa}{2c}$: in his scilicet locis maxime a recta $BAb$ remotis curvatura est maxima.

\originalsection{22}Etsi autem pro puncto $C$ constat abscissa $AE=c$, tamen distantia $EC$ nisi per integrationem æquationis
\[dy=\frac{(aa-cc+xx)dx}{\sqrt{(cc-xx)(2aa-cc+xx)}}\]
definiri non potest. Si enim post integrationem ponatur $x=c$; valor ipsius $y$ dabit distantiam $CE$, quæ bis sumpta præbebit distantiam $AB$, seu intervallum $Dd$, inter diametros interjacens. Simili modo
\sourcepage{259}integratione opus erit ad laminæ incurvatæ $AC$ longitudinem determinandam. Cum enim posito arcu $AM=s$, sit $ds=\frac{aa\,dx}{\sqrt{(cc-xx)(2aa-cc+xx)}}$, hujus integrale, posito $x=c$, dabit longitudinem curvæ $AC$.

\originalsection{23}Cum autem istæ formulæ integrationem non admittant, per approximationem valores intervalli $AD$ \& arcus curvæ $AC$ commode exprimere nitamur. Ponamus in hunc finem $\sqrt{cc-xx}=z$; eritque $PM=y=\int\frac{(aa-zz)dx}{z\sqrt{2aa-zz}}$, \& $AM=s=\int\frac{aa\,dx}{z\sqrt{2aa-zz}}$. Est vero per seriem
\[\frac1{\sqrt{2aa-zz}}=\frac1{a\sqrt2}\left(1+\frac14\times\frac{zz}{aa}+\frac{1.3}{4.8}\times\frac{z^4}{a^4}+\frac{1.3.5}{4.8.12}\times\frac{z^6}{a^6}+\&c.\right);\]
unde fiet
\[\begin{aligned}
s&=\frac1{\sqrt2}\int dx\left(\frac az+\frac14\times\frac za+\frac{1.3}{4.8}\times\frac{z^3}{a^3}+\frac{1.3.5}{4.8.12}\times\frac{z^5}{a^5}+\&c.\right),\\
s-y&=\frac1{\sqrt2}\int dx\left(\frac za+\frac14\times\frac{z^3}{a^3}+\frac{1.3}{4.8}\times\frac{z^5}{a^5}+\frac{1.3.5}{4.8.12}\times\frac{z^7}{a^7}+\&c.\right).
\end{aligned}\]

\originalsection{24}Quia autem hæc integralia tantum pro casu $x=c$ desideramus; quo casu fit $z=0$, ea commode ope peripheriæ Circuli exprimi poterunt. Posita enim ratione diametri ad peripheriam $=1:\pi$, erit $\int\frac{dx}z=\int\frac{dx}{\sqrt{cc-xx}}=\frac\pi2$; posito post integrationem $x=c$. Pari modo autem sequentia integralia ita determinabuntur, ut sit
\[\begin{aligned}
\int zdx&=\frac12\times\frac\pi2cc,\\
\int z^3dx&=\frac{1.3}{2.4}\times\frac\pi2c^4,\\
\int z^5dx&=\frac{1.3.5}{2.4.6}\times\frac\pi2c^6,\\
\int z^7dx&=\frac{1.3.5.7}{2.4.6.6}\times\frac\pi2c^8.\\
&\&c.
\end{aligned}\]
\sourcepage{260}His ergo integralibus in subsidium vocatis, erit:
\[\begin{aligned}
AC&=\frac{\pi a}{2\sqrt2}\left(1+\frac{1.1}{2.2}\times\frac{cc}{2aa}+\frac{1.1.3.3}{2.2.4.4}\times\frac{c^4}{4a^4}+\&c.\right),\\
AC-AD&=\frac{\pi a}{2\sqrt2}\left(\frac{cc}{2aa}+\frac{1.3}{2.4}\times\frac{c^4}{4a^4}+\frac{1.3.3.5}{2.4.4.6}\times\frac{c^6}{8a^6}+\&c.\right).
\end{aligned}\]
Ex his ergo reperiuntur $AD$ \& $AC$ ut sequitur:
\[\begin{aligned}
AC&=\frac{\pi a}{2\sqrt2}\left(1+\frac{1^2}{2^2}\times\frac{cc}{2aa}+\frac{1^2.3^2}{2^2.4^2}\times\frac{c^4}{4a^4}+\frac{1^2.3^2.5^2}{2^2.4^2.6^2}\times\frac{c^6}{8a^6}+\&c.\right),\\
AD&=\frac{\pi a}{2\sqrt2}\left(1-\frac{1^2}{2^2}\times\frac31\times\frac{cc}{2aa}-\frac{1^2.3^2}{2^2.4^2}\times\frac53\times\frac{c^4}{4a^4}\right.\\
&\qquad\left.-\frac{1^2.3^2.5^2}{2^2.4^2.6^2}\times\frac75\times\frac{c^6}{8a^6}-\&c.\right).
\end{aligned}\]
Si itaque detur $AE=c$, \& $AD=b$, ex his æquationibus \& recta constans $a$ \& longitudo curvæ $AC$ definietur. Vicissim autem ex data longitudine curvæ $AC$, \& recta $a$, per quam vis inflectens determinatur, reperiri poterunt rectæ $AD$ \& $CD$.

\originalsection{25}\marginnote{Species prima.}Quoniam igitur speciem primam constituimus, si in æquatione generali $dy=\frac{(aa-cc+xx)dx}{\sqrt{(cc-xx)(2aa-cc+xx)}}$ fuerit $c=0$, seu $\frac ac=\infty$, quo casu linea resultat repræsentans statum laminæ Elasticæ naturalem; ad eandem speciem primam referamus quoque eos casus, quibus $c$ est quantitas quamminima, ita ut præ $a$ pro evanescente haberi queat. Quia ergo $x$ ipsam $c$ superare nequit; etiam $x$ præ $a$ evanescet, ideoque ista prodibit æquatio $dy=\frac{a\,dx}{\sqrt{2(cc-xx)}}$, cujus integrale est $y=\frac a{\sqrt2}\operatorname{Asin}.\frac xc$, quæ est æquatio pro curva Trochoide in infinitum elongata. Fiet autem $AD=\frac{\pi a}{2\sqrt2}$, a qua ipsa curvæ longitudo infinite parum tantum discrepat, propterea quod angulus $DAM$ est infinite parvus. Sit longitudo laminæ $ACB=2f$, ejusque elasticitas absoluta $=Ekk$; ob $f=\frac{\pi a}{2\sqrt2}$, erit vis ad hanc curvaturam infinite parvam laminæ inducendam
\sourcepage{261}requisita finitæ magnitudinis \& quidem $=\frac{Ekk}{ff}\times\frac{\pi\pi}4$. Scilicet si extremitates $A$ \& $B$ colligentur filo $AB$, hoc filum contrahi debebit vi $=\frac{Ekk}{ff}\times\frac{\pi\pi}4$.

\originalsection{26}\marginnote{Species secunda.}Secundam speciem constituat casus, quo $c>0$, attamen $c<a$; scilicet si $c$ contineatur intra limites $0$ \& $a$. His enim casibus angulus $DAM$ recto erit minor; est namque anguli $PAM$ sinus, seu anguli $DAM$ cosinus $=\frac{aa-cc}{aa}$. Hoc ergo casu, forma lineæ curvæ talis fere erit qualem \hyperlink{addfig:6}{Figura 6}, repræsentat. Quia igitur est $c<a$ erit $\frac{cc}{2aa}<\frac12$; cum vero sit $\frac{cc}{2aa}>0$, erit utique $AC=f>\frac{\pi a}{2\sqrt2}$, unde $aa<\frac{8ff}{\pi\pi}$; quare vis, qua extremitates laminæ $A$ \& $B$ ope fili $AB$ ad se invicem attrahuntur, major erit quam casu præcedente, nempe $>\frac{Ekk}{ff}\times\frac{\pi\pi}4$.

\originalsection{27}\marginnote{Species tertia.}In tertia specie unicum complector casum, quo $c=a$, quia hoc casu axis $AP$ curvam in puncto $A$ tangit: hæcque species singulare nomen curvæ Elasticæ rectangulæ obtinuit. Erit ergo $dy=\frac{xx\,dx}{\sqrt{a^4-x^4}}$, \& $ds=\frac{aa\,dx}{\sqrt{a^4-x^4}}$; hoc igitur casu $AD$ \& $AC$ ita se habebunt ut sit:
\[\begin{aligned}
AC=f&=\frac{\pi a}{2\sqrt2}\left(1+\frac{1^2}{2^2}\times\frac12+\frac{1^2.3^2}{2^2.4^2}\times\frac14+\frac{1^2.3^2.5^2}{2^2.4^2.6^2}\times\frac18+\&c.\right),\\
AD=b&=\frac{\pi a}{2\sqrt2}\left(1-\frac{1^2}{2^2}\times\frac3{1.2}-\frac{1^2.3^2}{2^2.4^2}\times\frac5{3.4}-\frac{1^2.3^2.5^2}{2^2.4^2.6^2}\times\frac7{5.8}-\&c.\right).
\end{aligned}\]
Quanquam autem hinc, neque $b$, neque $f$ per $a$ accurate assignari potest; tamen alibi insignem relationem inter has quantitates locum habere demonstravi. Scilicet ostendi esse $4bf=\pi aa$, seu rectangulum ex $AD$ \& $AC$ formatum erit æquale areæ Circuli cujus diameter est $=AE$. Reperietur autem, calculum subducendo, proxime $f=\frac{5a}6\times\frac\pi2$, ita ut sit $a=\frac{12f}{5\pi}$; hinc vis \sourcepage{262}qua laminæ extremitates $A$, $B$ ad se invicem contrahi debent, erit $=\frac{Ekk}{ff}\times\frac{25}{72}\pi\pi$. Propius vero reperitur $f=\frac{\pi a}{2\sqrt2}.1{,}1803206$, hincque $b=\frac{\pi aa}{4f}=\frac a{\sqrt2}\times1{,}1803206$; unde in numeris puris erit $\frac fa=1{,}311006$, \& $\frac ba=0{,}834612$.

\originalsection{28}\marginnote{Species quarta.}Si $c>a$, orietur species quarta, eousque patens, quoad fiat $AD=b=0$; qui alter limes ipsius $c$ definietur per hanc æquationem:
\[1=\frac{1^2}{2^2}\times\frac31\times\frac{cc}{2aa}+\frac{1^2.3^2}{2^2.4^2}\times\frac53\times\frac{c^4}{4a^4}+\frac{1^2.3^2.5^2}{2^2.4^2.6^2}\times\frac75\times\frac{c^6}{8a^6}+\&c.\]
\marginnote{\hyperlink{addfig:7}{Fig.~7}}In hac ergo specie cum sit $c>a$; curva in $A$ supra axem $AE$ ascendet, angulumque constituet $PAM$, cujus sinus erit $=\frac{cc-aa}{aa}$; mox autem videbimus hunc angulum $PAM$ minorem esse quam $40^\circ$, $41'$; quoniam si hunc valorem acquirit, intervallum $AD$ evanescit, quem casum ad speciem quintam refero. Hinc in specie quarta continentur casus quibus $\frac{cc}{aa}$ inter hos limites $1$ \& $1{,}651868$ comprehenditur. Harum autem curvarum forma ex figura intelligitur; dummodo notetur, quo propius $\frac{cc}{aa}$ ad posteriorem limitem $1{,}651868$ accesserit; eo minus esse futurum intervallum $AD$, eoque propius laminæ terminos $A$ \& $B$ ad se invicem adduci. Fieri ergo potest ut laminæ gibbositates $m$ \& $R$, item $M$ \& $r$, se mutuo non solum tangant, sed etiam intersecent, atque hujusmodi intersectiones in infinitum multiplicabuntur, donec omnes diametri $DC$, $dc$ coincidant, atque cum axe $AE$ confundantur.

\originalsection{29}\marginnote{Species quinta. \hyperlink{addfig:8}{Fig.~8}}Hoc si evenerit, orietur species quinta, cujus natura hac exprimetur æquatione inter coordinatas $AP=x$ \& $PM=y$;
\[dy=\frac{(cc-aa-xx)dx}{\sqrt{(cc-xx)(2aa-cc+xx)}},\]
existente hac inter $a$ \& $c$ relatione, ut sit intervallum $AD=b=0$. Ponatur
\sourcepage{263}$\frac{cc}{2aa}=v$, atque $v$ ex hac æquatione infinita definiri debet
\[1=\frac{1.3}{2.2}v+\frac{1.1.3.5}{2.2.4.4}v^2+\frac{1.1.3.3.5.7}{2.2.4.4.6.6}v^3+\&c.\]
Quærantur primum per methodos cuique solitas, vel saltem tentando, limites inter quos verus valor ipsius $v$ contineatur, atque hujusmodi limites reperientur $v=0{,}824$ \& $v=0{,}828$. Quod si jam uterque substituatur in æquatione ex erroribus binis oriundis, concludetur tandem fore $v=0{,}825934=\frac{cc}{2aa}$; unde fit $\frac{cc}{aa}=1{,}651868$ \& $\frac{cc-aa}{aa}=0{,}651868$; quæ expressio cum sit sinus anguli $PAM$, ex Tabulis reperietur hic angulus $=40^\circ$, $41'$; ideoque hujus duplum, seu angulus $MAN$, erit $=81^\circ$, $82'$. Quare si laminæ elasticæ extremitates eousque ad se invicem adducantur, ut se contingant; tum curvam $AMCNA$ formabunt, \& ambæ extremitates in $A$ angulum constituent $=81^\circ$, $22'$.

\originalsection{30}\marginnote{Species sexta. \hyperlink{addfig:9}{Fig.~9}}Si ambæ extremitates laminæ $A$ \& $B$, postquam ad se invicem fuerint adductæ, aucta vi in plagas contrarias a se invicem diducantur; orietur curva hujus formæ $AMCNB$, quæ speciem sextam constituat. In curvis ergo ad hanc speciem pertinentibus, erit $\frac{cc}{2aa}>0{,}825934$; ita tamen ut sit $\frac{cc}{2aa}<1$. Quod si enim sit $cc=2aa$ orietur species septima mox explicanda. Erit ergo in his curvis angulus $PAM$, quem curva in $A$ cum axe constituit major quam $40^\circ$, $41'$, minor tamen recto: cum enim ejus sinus sit $=\frac{cc-aa}{aa}$, ob $cc<2aa$, sinus iste necessario est minor sinu toto; neque ergo angulus $PAM$ rectus fieri potest, nisi ponatur $cc=2aa$.

\originalsection{31}\marginnote{Species septima. \hyperlink{addfig:10}{Fig.~10}}Sit jam $cc=2aa$, quo casu species septima constituitur, atque natura curvæ exprimetur hac æquatione
\[dy=\frac{(aa-xx)dx}{x\sqrt{2aa-xx}};\]
ex qua colligitur, curvæ ramos $A$ \& $B$ infinitum extendi ita, ut recta $AB$ fiat curvæ asymtota. Fiet ergo uterque ramus $AMC$ \& $BNC$ infinitus, id quod ex \sourcepage{264}serie supra pro arcu $AC$ inventa intelligitur; erit enim
\[AC=\frac{\pi a}{2\sqrt2}\left(1+\frac{1^2}{2^2}+\frac{1^2.3^2}{2^2.4^2}+\frac{1^2.3^2.5^2}{2^2.4^2.6^2}+\&c.\right),\]
cujus seriei summa est infinita. Quod si igitur laminæ longitudo $AC$ fuerit finita $=f$, necesse est ut sit $a=0$, hincque etiam $CD=c=0$; lamina ergo, postquam in nodum fuerit incurvata, hoc casu iterum in directum extendetur, ad quam extensionem opus erit vi infinita. Sin autem lamina fuerit infinite longa, curvam formabit nodatam ad asymtotam $AB$ convergentem, existente $CD=c$. Æquatio autem pro hac curva ope logarithmorum integrari potest, obtinebitur enim
\[y=\sqrt{cc-xx}-\frac c2l\frac{c+\sqrt{cc-xx}}x,\]
sumptis abscissis $x$ in ipsa diametro $DC$; ita ut sit $DQ=x$, \& $QM=y$; evanescit enim applicata $y$, posito $x=CD=c$. In nodo autem $O$ applicata $y$ pariter evanescit: ad quem locum inveniendum, ponatur $\frac{2\sqrt{cc-xx}}c=l\frac{c+\sqrt{cc-xx}}x$. Sit $\varphi$ angulus cujus cosinus $=\frac xc$ \& sinus $=\frac{\sqrt{cc-xx}}c$, erit $2\sin.\varphi=l\operatorname{tang}.(45^\circ+\frac12\varphi)$, qui logarithmus ex hyperbolicorum genere sumi debet; cujusmodi Canon si deficiat, sumatur ex Canone vulgari logarithmus tangentis anguli $45^\circ+\frac12\varphi$, a cujus characteristica denarius auferatur, sitque residuum $=\omega$; quo facto erit $2\sin.\varphi=\omega.2{,}30258509$: sumendis ergo iterum logarithmis vulgaribus, erit $l2+l\sin.\varphi=l\omega+0{,}3622156886$, seu $l\sin.\varphi=l\omega+0{,}0611856930$. Hoc artificio tentando, mox vero proximus valor anguli $\varphi$ elicietur; unde porro per regulam falsi verus valor anguli $\varphi$, ex eoque abscissa $x=DO$ definietur. Reperitur autem hoc modo angulus $\varphi=73^\circ$, $14'$, $12''$, unde prodit $\frac xc=0{,}2884191$, \& $\frac{\sqrt{cc-xx}}c=0{,}9575042$; angulus vero $QOM$ fit $=2\varphi-90^\circ=56^\circ$, $28'$, $24''$, ideoque angulus $MON=112^\circ$, $56'$, $48''$. Cum igitur specie quinta angulus nodi esset $81^\circ$, $22'$, in
\sourcepage{265}specie sexta angulus nodi $MON$ continebitur inter limites $81^\circ$, $22'$ \& $112^\circ$, $56'$, $48''$. In specie quarta autem siquidem detur nodus, erit ejus angulus minor quam $81^\circ$, $22'$.

\originalsection{32}\marginnote{Species octava. \hyperlink{addfig:11}{Fig.~11}}Sit jam $cc>2aa$, puta $cc=2aa+gg$; erit æquatio pro curva, ob $aa=\frac{cc-gg}2$
\[dy=\frac{(xx-\frac12cc-\frac12gg)dx}{\sqrt{(cc-xx)(xx-gg)}},\]
qua æquatione species octava continetur, eritque, si recta $dDd$ repræsentet directionem vis sollicitantis, $x=DQ$ \& $y=QM$. Primum ergo patet applicatam $y$ realem esse non posse, nisi sit $x>g$, tum vero $x$ non potest excedere rectam $DC=c$, unde sumpta $DF=g$, tota curva continebitur inter rectas ipsi $dd$ parallelas per puncta $C$ \& $F$ ductas, quæ curvam simul tangent. Perinde autem est utra rectarum $c$ \& $g$ sit major, dummodo fuerint inæquales, æquatio enim non variatur si rectæ $c$ \& $g$ inter se permutentur. Deinde vero hæc curva quoque habebit infinitas diametros inter se parallelas $DC$, $dc$, $dc$, \& quæ per singula puncta $G$ \& $H$ ducuntur rectæ pariter ad $dDd$ normales; nusquam autem per totam curvam dabitur punctum flexus contrarii, ideoque continua curvatura utrinque in infinitum progreditur, uti figura indicat; anguli autem qui in nodis constituuntur $MON$ majores erunt quam $112^\circ$, $56'$, $42''$.

\originalsection{33}\marginnote{Species nona.}Cum in hac specie non solum contineantur casus quibus $gg<cc$, sed etiam quibus $gg>cc$, unicus adhuc superest casus quo $c=g$, quo quidem tota curva in spatium evanescens, ob $CF=0$, redigitur. Quod si autem utramque $c$ \& $g$ statuamus infinitam, ita tamen ut earum differentia fiat finita, curva finitum spatium occupabit. Ad eam ergo inveniendam, ponatur $g=c-2h$, \& $x=c-h-t$, atque ob $c=\infty$, quantitates $h$ \& $t$ vero finitæ, erit $\frac12cc+\frac12gg=cc-2ch$; \& $xx-\frac12cc-\frac12gg=-2ct$; tum vero $cc-xx=2c(h+t)$, \& $xx-gg=2c(h-t)$; ex quibus sequens prodibit æquatio, $dy=\frac{t\,dt}{\sqrt{hh-tt}}$, pro Circulo. Lamina \sourcepage{266}elastica ergo hoc casu in Circulum incurvatur, uti supra jam annotavimus; Circulus ergo speciem nonam atque ultimam constituet.

\originalsection{34}His enumeratis speciebus, facile erit pro quovis casu oblato assignare, ad quamnam speciem curva formata pertineat. \marginnote{\hyperlink{addfig:12}{Fig.~12}}Sit lamina elastica in $G$ muro infixa, termino vero $A$ appendatur pondus $P$, quo lamina in figuram $GA$ incurvetur. Ducatur tangens $AT$, atque ex angulo $TAP$ totum judicium erit petendum. Si enim hic angulus fuerit acutus, referetur curva ad speciem secundam; sin sit rectus ad tertiam, eritque elastica rectangula. Quod si angulus $TAP$ fuerit obtusus, minor tamen quam $130^\circ$, $41'$, curva ad speciem quartam pertinebit; ad quintam autem si angulus $TAP$ sit $=130^\circ$, $41'$; sin autem angulus $TAP$ major fuerit, curva sub specie sexta continebitur. Ad septimam vero pertineret, si iste angulus fieret duobus rectis æqualis; quod autem nunquam fieri potest. Hæc igitur species cum duabus sequentibus produci nequit laminæ immediate pondus appendendo.

\originalsection{35}\marginnote{\hyperlink{addfig:3}{Fig.~3}}Ut igitur pateat quomodo reliquæ species laminam incurvando produci queant, laminæ in $B$ fixæ, non immediate, sed virgæ rigidæ $AC$ cum laminæ termino $A$ firmissime connexæ in $C$ appendatur pondus $P$, quod trahat in directione $CD$. Sit intervallum $AC=h$, elasticitas laminæ absoluta $=Ekk$, \& anguli $MAP$ quem lamina in $A$ cum horizontali constituit, sinus $=m$. His positis, si ponatur abscissa $AP=t$ \& applicata $PM=y$, reperietur pro curva ista æquatio
\[dy=\frac{dt(mEkk-Pht-\frac12Ptt)}{\sqrt{E^2k^4-(mEkk-Pht-\frac12Ptt)^2}}.\]
Ponatur jam $CP=x=h+t$, quo æquatio ad formam qua in divisione specierum usi sumus, reducatur; erit
\[dy=\frac{dx(mEkk+\frac12Phh-\frac12Pxx)}{\sqrt{E^2k^4-(mEkk+\frac12Phh-\frac12Pxx)^2}},\]
quæ comparata cum forma
\[dy=\frac{dx(aa-cc+xx)}{\sqrt{(cc-xx)(2aa-cc+xx)}},\]
\sourcepage{267}seu $dy=\frac{dx(aa-cc+xx)}{\sqrt{a^4-(aa-cc+xx)^2}}$ dabit $\frac12Paa=Ekk$, seu $aa=\frac{2Ekk}P$; \& $\frac12Pcc-\frac12Paa=mEkk+\frac12Phh$; ergo $cc=\frac{2(1+m)Ekk}P+hh$.

\originalsection{36}Curva ergo ad speciem secundam pertinebit, si fuerit $\frac{2mEkk}P+hh<0$; seu $P<\frac{-2mEkk}{hh}$: nisi ergo angulus $PAM$ sit negativus, vis $P$ negativa esse atque virga in $C$ sursum trahi debet. Ad speciem tertiam curvatura pertinebit, si $P=\frac{-2mEkk}{hh}$. Quarta autem species prodibit si fuerit $2mEkk+Phh>0$, simul vero $2mEkk+Phh<2\alpha Ekk$, existente $\alpha=0{,}651868$. Sin autem sit $P=\frac{2(\alpha-m)Ekk}{hh}$, tum curva ad speciem quintam pertinebit. Quod si vero fuerit $Phh>2(\alpha-m)Ekk$, simul vero $Phh<2(1-m)Ekk$, curva ad speciem sextam est referenda. Septimaque species proveniet, si $Phh=2(1-m)Ekk$. Octava autem species obtinebitur, si $Phh>2(1-m)Ekk$; quare si angulus $PAM$ fuerit rectus, ob $1-m=0$, curva semper ad speciem octavam pertinebit. Species denique nona orietur, si fuerit $h=\infty$; uti jam supra annotavi.

\originalsection{37}\marginnote{\hyperlink{addfig:13}{Fig.~13}}Quæ ante de specie prima sunt annotata inservire possunt viribus columnarum dijudicandis. Sit enim $AB$ columna super basi $A$ verticaliter posita, gestans pondus $P$. Quod si jam columna ita sit constituta ut prolabi nequeat; ab onere $P$, si fuerit nimis magnum, nil aliud erit metuendum, nisi columnæ incurvatio; hoc ergo casu columna spectari poterit tanquam elasticitate prædita. Sit igitur elasticitas absoluta columnæ $=Ekk$: ejusque altitudo $AB=2f=a$, atque supra §.~25 vidimus, vim requisitam ad hanc columnam vel minimum inclinandam esse $=\frac{\pi\pi Ekk}{4ff}=\frac{\pi\pi}{aa}.Ekk$. Nisi ergo onus gestandum $P$ majus \sourcepage{268}sit quam $E.\frac{\pi\pi kk}{aa}$, nulla prorsus incurvatio erit metuenda; contra vero si pondus $P$ fuerit majus, columna incurvationi resistere non poterit. Manente autem elasticitate columnæ, atque adeo ejus crassitie eadem; pondus $P$, quod sine periculo gestare valet, erit reciproce ut quadratum altitudinis columnæ: columnaque duplo altior quartam tantum oneris partem gestare poterit. Hæc igitur præcipue in usum vocari possunt circa columnas ligneas, quippe quæ incurvationi sunt obnoxiae.

\originalsection{38}\marginnote{\hyperlink{addfig:14}{Fig.~14}}Quo autem vis atque incurvatio cujusque laminæ elasticæ a priori determinari queat; necesse est ut elasticitas absoluta, quam hactenus per $Ekk$ expressimus, sit cognita; id quod unico experimento commode præstabitur. Infigatur lamina elastica uniformis $FH$, cujus elasticitatem absolutam investigari oportet, altero termino $F$ parieti firmo $GK$; ita ut situm teneat horizontalem $FH$; hic enim gravitatem naturalem negligere liceat. Alteri termino $H$ appendatur pondus pro arbitrio sumptum $P$, quo lamina in statum $AF$ incurvetur. Sit longitudo laminæ $AF=HF=f$, recta horizontalis $AG=g$, \& verticalis $GF=h$; qui valores omnes per experimentum erunt cogniti. Comparetur jam hæc curva cum æquatione generali
\[dy=\frac{(cc-aa-xx)dx}{\sqrt{(cc-xx)(2aa-cc+xx)}},\]
in qua si fuerint $a$ \& $c$ per $f$, $g$, $h$, definitæ; erit vis incurvans $P=\frac{2Ekk}{aa}$; ideoque elasticitas absoluta $Ekk=\frac12Paa$.

\originalsection{39}Quia jam tangens in $F$ est horizontalis, erit hic $\frac{dy}{dx}=0$; ideoque $x=\sqrt{cc-aa}$. Hinc ergo erit $AG=g=\sqrt{cc-aa}$, \& $aa=cc-gg$; ideoque
\[dy=\frac{(gg-xx)dx}{\sqrt{(cc-xx)(cc-2gg+xx)}},\]
posito autem hic $x=g$, fieri debebit $y=GF=h$, seu $s=AF=f$; est vero $ds=\frac{(cc-gg)dx}{\sqrt{(cc-xx)(cc-2gg+xx)}}$. Jam si pondus $P$ sumatur valde parvum, ut lamina paulisper
\sourcepage{269}tantum deprimatur; tum erit $c$ quantitas valde magna, ideoque erit proxime
\[\begin{aligned}
\frac1{\sqrt{(cc-xx)(cc-2gg+xx)}}&=(c^4-2ccgg+2ggxx-x^4)^{-\frac12}\\
&=\frac1{cc}+\frac{gg}{c^4}-\frac{ggxx}{c^6}+\frac{x^4}{2c^6},
\end{aligned}\]
ideoque integrando quoque proxime:
\[\begin{aligned}
s&=\frac{(cc-gg)x}{cc}+\frac{(cc-gg)ggx}{c^4}-\frac{(cc-gg)ggx^3}{3c^6}+\frac{(cc-gg)x^5}{10c^6},\\
y&=\frac{ggx}{cc}+\frac{g^4x}{c^4}-\frac{g^4x^3}{3c^6}+\frac{ggx^5}{10c^6}\\
&\qquad-\frac{x^3}{3cc}-\frac{ggx^3}{3c^4}+\frac{ggx^5}{5c^6}-\frac{x^7}{14c^6}.
\end{aligned}\]
Sit nunc $x=g$, fietque $f=g-\frac{37g^5}{30c^4}$ \& $h=\frac{2g^3}{3cc}+\frac{2g^5}{3c^4}$. Quod si ergo recta $FG=h$ in usum vocetur, erit $cc=\frac{2g^3}{3h}$, \& $aa=\frac{g(2gg-3gh)}{3h}$: unde elicitur elasticitas absoluta
\[Ekk=\frac{Pgg(2g-3h)}{6h};\]
qui valor a vero vix sensibiliter discrepabit, dummodo laminæ curvatura non nimis magna inducatur.

\originalsection{40}Hæc autem elasticitas absoluta $Ekk$ primum pendet ab natura materiæ, ex qua lamina est fabrefacta; unde alia materia magis, alia minus elatere prædita dici solet. Secundo quoque ita pendet a laminæ latitudine, ut expressio $Ekk$ ubique latitudini laminæ debeat esse proportionalis, si cetera sint paria. Tertio verum crassities laminæ plurimum confert ad valorem ipsius $Ekk$ determinandum, quæ ita comparata esse videtur, ut, ceteris paribus, $Ekk$ sit ut crassitiei quadratum. Conjunctim ergo tenebit expressio $Ekk$ rationem compositam ex ratione elateris materiæ, latitudinis laminæ simplici, ac duplicata crassitiei laminæ. Hinc per experimenta quibus latitudinem \& crassitiem \sourcepage{270}metiri licet, omnium materiarum elasticitates inter se comparari ac determinari poterunt.

\originalsection{41}\marginnote{\hyperlink{addfig:2}{Fig.~2}}Quemadmodum igitur hactenus laminæ, cujus curvaturam determinavi, elasticitatem absolutam $Ekk$ per totam longitudinem constantem posui; ita solutio eadem methodo poterit absolvi, si quantitas $Ekk$ utcunque ponatur variabilis. Scilicet si elasticitas absoluta fuerit ut functio quæcunque portionis laminæ $AM$, quæ functio sit $=S$, posito arcu $AM=s$; atque existente radio osculi in $M=R$; curva $AM$, quam lamina induit, ita erit comparata, ut in ea, inter omnes alias ejusdem longitudinis, sit $\int\frac{Sds}{RR}$ minimum. Solvetur ergo iste casus per formulam secundam generalem. Sit $dy=pdx$; $dp=qdx$; at $dS=Tds$, atque, inter omnes curvas in quibus est $\int dx\sqrt{1+pp}$ ejusdem magnitudinis, ea determinari debebit, in qua sit $\int\frac{Sqqdx}{(1+pp)^{5:2}}$ minimum. Prior formula $\int dx\sqrt{1+pp}$ dat pro formula differentiali $\frac1{dx}d.\frac p{\sqrt{1+pp}}$. Altera vero $\int\frac{Sqqdx}{(1+pp)^{5:2}}$ cum $\int Zdx$ comparata dabit $Z=\frac{Sqq}{(1+pp)^{5:2}}$. Cum igitur positum sit $dZ=Ld\pi+Mdx+Ndy+Pdp+Qdq$, $\pi=\int[Z]dx$, \& $d[Z]=[M]dx+[N]dy+[P]dp$; erit $Ld\pi=\frac{qqTds}{(1+pp)^{5:2}}$; unde $L=\frac{qqT}{(1+pp)^{5:2}}$; $d\pi=ds=dx\sqrt{1+pp}$; ideoque $[Z]=\sqrt{1+pp}$; $[M]=0$, $[N]=0$, \& $[P]=\frac p{\sqrt{1+pp}}$. Deinde vero est $M=0$, $N=0$, $P=-\frac{5Sqqpdp}{(1+pp)^{7:2}}$, \& $Q=\frac{2Sq}{(1+pp)^{5:2}}$; ita ut sit $dZ=\frac{qqdS}{(1+pp)^{5:2}}+Pdp+Qdq$.

\originalsection{42}Jam sumatur integrale
\[\int Ldx=\int\frac{qqTdx}{(1+pp)^{5:2}}=\int\frac{qqdS}{(1+pp)^3};\]
\sourcepage{271}sitque $H$ ejus valor, si ponatur $x=a$, cujus quidem constantis $a$ consideratio mox ex calculo rursus evanescet. Erit ergo $V=H-\int\frac{qqdS}{(1+pp)^3}$. Unde valor differentialis fiet $=-\frac{dP}{dx}-\frac1{dx}d.[P]V+\frac{ddQ}{dx^2}$. Quamobrem ex his duobus valoribus differentialibus nascetur hæc æquatio pro curva quæsita
\[\frac\alpha{dx}d.\frac p{\sqrt{1+pp}}=+\frac{dP}{dx}+\frac1{dx}d.[P]V-\frac{ddQ}{dx^2},\]
quæ integrata dat
\[\frac{\alpha p}{\sqrt{1+pp}}+\beta=P+[P]V-\frac{dQ}{dx},\]
sive
\[\frac{\alpha p}{\sqrt{1+pp}}+\beta=\frac{Hp}{\sqrt{1+pp}}-\frac p{\sqrt{1+pp}}\int\frac{qqdS}{(1+pp)^3}+P-\frac{dQ}{dx},\]
ubi constans $H$ alias determinata in constante arbitraria $\alpha$ comprehendi potest, quo ipso constans $a$ ex calculo egreditur. Idcirco ergo prodibit hæc æquatio
\[\frac{\alpha p}{\sqrt{1+pp}}+\beta=P-\frac{dQ}{dx}-\frac p{\sqrt{1+pp}}\int\frac{qqdS}{(1+pp)^3}.\]

\originalsection{43}Multiplicetur hæc æquatio per $dp=qdx$, atque prodibit:
\[\frac{\alpha p\,dp}{\sqrt{1+pp}}+\beta dp=Pdp-Qdq-\frac{p\,dp}{\sqrt{1+pp}}\int\frac{qqdS}{(1+pp)^3}.\]
Cum autem sit $dZ=\frac{qqdS}{(1+pp)^{5:2}}+Pdp+Qdq$, erit $Pdp=dZ-Qdq-\frac{qqdS}{(1+pp)^{5:2}}$; quo valore substituto, emerget æquatio integrabilis hæc:
\[\begin{aligned}
\frac{\alpha p\,dp}{\sqrt{1+pp}}+\beta dp&=dZ-qdQ-Qdq-\frac{qqdS}{(1+pp)^{5:2}}\\
&\qquad-\frac{p\,dp}{\sqrt{1+pp}}\int\frac{qqdS}{(1+pp)^3},
\end{aligned}\]
cujus integralis est:
\[\alpha\sqrt{1+pp}+\beta p+\gamma=Z-Qq-\sqrt{1+pp}\int\frac{qqdS}{(1+pp)^3},\]
seu
\sourcepage{272}
\[\alpha\sqrt{1+pp}+\beta p+\gamma=\frac{-Sqq}{(1+pp)^{5:2}}-\sqrt{1+pp}\int\frac{qqdS}{(1+pp)^3}.\]
Quo signum integrale tollamus, divisa æquatione per $\sqrt{1+pp}$, ea denuo differentietur;
\[\frac{\beta dp}{(1+pp)^{3:2}}-\frac{\gamma p\,dp}{(1+pp)^{3:2}}+\frac{2qqdS}{(1+pp)^3}+\frac{2Sqdq}{(1+pp)^3}-\frac{6Spqqdp}{(1+pp)^4}=0,\]
quæ per $\frac{(1+pp)^{3:2}}{2q}$ multiplicata, præbet:
\[\frac{\beta dp}{2q}-\frac{\gamma p\,dp}{2q}+\frac{qdS+Sdq}{(1+pp)^{3:2}}-\frac{3Spqdp}{(1+pp)^{5:2}}=0;\]
cujus, ob $dp=qdx$ \& $dy=pdx$, integrale erit
\[\alpha+\frac12\beta x-\frac12\gamma y+\frac{Sq}{(1+pp)^{3:2}}=0.\]
At est $\frac{-(1+pp)^{3:2}}q=\text{radio osculi }R$; unde constantes $\beta$ \& $\gamma$ duplicando, orietur hæc æquatio $\frac SR=\alpha+\beta x-\gamma y$: quæ æquatio apprime congruit cum ea, quam altera Methodus directa suppeditat. Exprimet enim $\alpha+\beta x-\gamma y$ momentum potentiæ incurvantis, recta quacunque pro axe assumpta, cui momento utique æqualis esse debet elasticitas absoluta $S$ per radium osculi $R$ divisa. Sic igitur non solum Celeb. Bernoulli observata proprietas Elasticæ plenissime est evicta; sed etiam formularum mearum difficiliorum usus summus in hoc Exemplo est declaratus.

\originalsection{44}\marginnote{\hyperlink{addfig:3}{Fig.~3}}Si ergo curva fuerit data, quam lamina inæquabiliter elastica a potentia $CD=P$, sollicitata format; hinc elasticitas absoluta laminæ in quovis loco poterit cognosci. Sumpta enim recta $CP$, quæ ad directionem vis sollicitantis est normalis, pro axe, ac posita $CP=x$, $PM=y$, arcu curvæ $AN=s$, \& radio osculi in $M=R$, ob momentum potentiæ $P$ ad punctum $M$ relatum $=Px$; erit $\frac SR=Px$; ideoque elasticitas absoluta in puncto $M$, quæ est $S$, $=PRx$. Hinc,
\sourcepage{273}cum, data curva, in singulis punctis detur radius osculi $R$, elasticitas absoluta in quovis loco innotescit. Quod si ergo materia laminæ, una cum crassitie, ubique fuerit eadem; latitudo autem sit variabilis: quia elasticitas absoluta latitudini est proportionalis, ex curva formata latitudo laminæ in singulis locis colligitur.

\originalsection{45}\marginnote{\hyperlink{addfig:15}{Fig.~15}}Sit ex lamina elastica excissa lingula triangularis $fAf$, ubique ejusdem crassitiei. Quoniam ergo latitudo $mm$, in quovis loco $M$, est longitudini $AM$ proportionalis; posita $AM=s$, erit elasticitas absoluta in $M$ ut $s$. Sit ea $=Eks$; atque laminæ termino $ff$ muro horizontaliter infixo, appendatur cuspidi $A$ pondus $P$; quo laminæ recta media $AF$ in curvam $FmA$ incurvetur, cujus natura quæritur. \marginnote{\hyperlink{addfig:14}{Fig.~14}}Positis autem in axe horizontali abscissa $Ap=x$, applicata $pm=y$, \& arcu $Am=s$; erit $Px=\frac{Eks}R$, denotante $R$ radium osculi in $m$. Multiplicetur hæc æquatio per $dx$, \& ob $R=\frac{-ds^3}{dx\,ddy}$, posito $dx$ constante, erit $Pxdx=\frac{-Eks\,dx^2ddy}{ds^3}$, seu $\frac{Pxdx}{Ek}+\frac{s\,dx^2ddy}{ds^3}=0$. At, cum sit
\[d.\frac{sdy}{ds}=\frac{sddy}{ds}-\frac{sdydds}{ds^2}+dy=\frac{s\,dx^2ddy}{ds^3}+dy;\]
ob $dds=\frac{dyddy}{ds}$, erit $\int\frac{s\,dx^2ddy}{ds^3}=\frac{sdy}{ds}-y$; unde integrando habebitur $\frac{Pxx}{2Ek}+a=\frac{-sdy}{ds}+y$.

\originalsection{46}Sit $dy=pdx$, erit $ds=dx\sqrt{1+pp}$, \& posito $\frac{2Ek}P=c$, fiet $a+\frac{xx}c=y-\frac{sp}{\sqrt{1+pp}}$; ideoque erit
\[\frac{a\sqrt{1+pp}}p+\frac{xx\sqrt{1+pp}}{cp}=\frac{y\sqrt{1+pp}}p-s;\]
quæ differentiata dat
\[\begin{aligned}
&\frac{-a\,dp}{pp\sqrt{1+pp}}+\frac{2xdx\sqrt{1+pp}}{cp}-\frac{xxdp}{cpp\sqrt{1+pp}}\\
&\quad=\frac{dy\sqrt{1+pp}}p-\frac{y\,dp}{pp\sqrt{1+pp}}-dx\sqrt{1+pp}=\frac{-y\,dp}{pp\sqrt{1+pp}}.
\end{aligned}\]
Hinc oritur $a-y=\frac{2pxdx(1+pp)}{c\,dp}-\frac{xx}c$. Ponatur $dp$ constans, \sourcepage{274}ac differentiando erit
\[-pdx=\frac{2pxddx(1+pp)}{c\,dp}+\frac{2pdx^2(1+pp)}{c\,dp}+\frac{2xdx(1+3pp)}c-\frac{2xdx}c,\]
seu
\[0=c\,dxdp+2xddx(1+pp)+2dx^2(1+pp)+6pxdx;\]
cujus æquationis autem resolutio ulterior non constat. Simplicissima autem pro curva est æquatio hæc $\frac{yds-sdy}{ds}=\frac{Pxx}{2Ek}$; quia enim posito $x=0$, \& $y$ \& $s$ evanescere debent, constans $a$ debet esse $=0$.

\originalsection{47}Hoc igitur modo curvatura laminæ, sive æqualiter sive inæqualiter elasticæ, determinatur, si ab una potentia sollicitetur; atque, quod præcipue est notandum, si lamina naturaliter fuerit in directum extensa. Quod si enim lamina in statu naturali jam fuerit curva; tum utique a vi sollicitante aliam curvaturam induet; ad quam inveniendam, præter sollicitationem atque elasticitatem, simul figuram ejus naturalem nosse oportet. Sit igitur lamina elastica naturaliter curva $Bma$, cujus quidem elasticitas sit ubique eadem, $=Ekk$; quæ a vi sollicitante $P$ in figuram $BMA$ incurvetur. \marginnote{\hyperlink{addfig:16}{Fig.~16}}Per $A$ ducatur recta $CAP$ ad directionem vis sollicitantis normalis, quæ habeatur pro axe; sitque intervallum $AC=c$, abscissa $AP=x$, applicata $PM=s$; erit momentum vis sollicitantis pro puncto $M=P(c+x)$.

\originalsection{48}Sit porro radius osculi curvæ quæsitæ in $M=R$; sumatur in statu naturali arcus $am=AM=s$, sitque in puncto $m$ radius osculi $=r$; qui ob curvam $amB$ cognitam dabitur per arcum $s$. In $M$ ergo, quia curvatura major est, radius osculi $R$ minor est quam $r$, atque excessus anguli elementaris in $M$ supra angulum in statu naturali erit $=\frac{ds}R-\frac{ds}r$, qui excessus erit effectus a potentia sollicitante productus. Quamobrem erit $P(c+x)=Ekk(\frac1R-\frac1r)$; quæ cum $r$ per $s$ detur, erit æquatio pro curva quæsita; quæ autem sic in genere spectata ulterius reduci non potest.

\originalsection{49}\sourcepage{275}Ponamus ergo laminam in statu naturali $amB$ habere figuram circularem; erit $r$ radius ejus circuli qui sit $=a$, unde fit $P(c+x)=Ekk(\frac1R-\frac1a)$. Multiplicetur hæc æquatio per $dx$, \& integretur; orietur $\frac P{Ekk}(\frac12xx+cx+f)=\frac{-dy}{ds}-\frac xa$: quæ æquatio, si loco $c$ scribatur $c+\frac{Ekka}P$; abibit in $\frac P{Ekk}(\frac12xx+cx+f)=\frac{-dy}{ds}$: quæ est eadem æquatio quam supra pro lamina naturaliter recta invenimus. Lamina ergo naturaliter circularis in easdem curvas incurvatur, quæ laminæ naturaliter rectæ inducuntur: tantum scilicet locus applicationis potentiæ, seu intervallum $AC=c$, pro utroque casu secundum datam legem variari debebit. Eædem ergo novem species curvarum prodibunt pro figuris, quas lamina naturaliter circularis inducere potest, quas supra numeravimus. Lamina enim circularis, si intervallum $AC$ capiatur infinitum, primum in lineam rectam extendi potest; tum quæcunque potentia insuper applicata eundem præstabit effectum, ac si sola laminæ elasticæ naturaliter rectæ applicaretur.

\originalsection{50}Ponamus autem, quæcunque sit laminæ figura naturalis, punctum $C$ infinite distare; ita ut momentum vis sollicitantis ubique sit idem, quod per $Ekk$ divisum ponatur $=\frac1b$; eritque $\frac1b=\frac1R-\frac1r$ \& $\frac1R=\frac1b+\frac1r$. Hinc fiet $\int\frac{ds}R=\frac sb+\int\frac{ds}r=\text{amplitudini arcus }AM$; sicuti $\int\frac{ds}r$ exprimit amplitudinem arcus $am$; quemadmodum quidem Celeb. \textit{Joh. Bernoulli} hoc \textit{amplitudinis} nomine in eximio Tractatu \textit{De motu reptorio} uti est solitus. Sit igitur $\frac sb+\int\frac{ds}r$ arcus in circulo, cujus radius $=1$ sumptus, qui ob $r$ per $s$ datum, quoque in $s$ erit cognitus. Hinc autem reperientur coordinatæ orthogonales $x$ \& $y$, ita ut sit
\sourcepage{276}
\[x=\int ds\sin.\left(\frac sb+\int\frac{ds}r\right),\quad\&\quad y=\int ds\cos.\left(\frac sb+\int\frac{ds}r\right);\]
unde curva quæsita per quadraturas construi poterit.

\originalsection{51}\marginnote{\hyperlink{addfig:17}{Fig.~17}}Hinc determinari potest figura $amB$, quam lamina in situ naturali habere debet, ut a potentia $P$ in directione $AP$ sollicitante in lineam rectam $AMB$ explicetur. Sumpta enim longitudine $AM=s$; erit momentum potentiæ sollicitantis pro puncto $M=Ps$; radius osculi autem in $M$, per hypothesin, erit infinitus seu $\frac1R=0$. Sumpto jam in statu naturali arcu $am=s$, positoque radio osculi in $m=r$; quia hæc curva convexitate sua rectam $AB$ spectat, in calculo præcedente poni debet $r$ negativum. Hinc erit $Ps=\frac{Ekk}r$, seu $rs=aa$; quæ est æquatio naturam curvæ $amB$ complectens.

\originalsection{52}Cum igitur sit $\frac1r=\frac s{aa}$; erit $\int\frac{ds}r=\frac{ss}{2aa}$; seu erit amplitudo arcus $am$ ut quadratum ipsius arcus. Hinc coordinatæ orthogonales $x$ \& $y$ pro hac curva $amB$ ita definientur, ut sit $x=\int ds\sin.\frac{ss}{2aa}$, \& $y=\int ds\cos.\frac{ss}{2aa}$: Scilicet in circulo, cujus radius $=1$, abscindi debet arcus $=\frac{ss}{2aa}$, cujus sinus \& cosinus ad coordinatas determinandas assumi debent. Ex eo autem quod radius osculi continuo decrescit, quo major capiatur arcus $am=s$, manifestum est curvam in infinitum non protendi, etiamsi arcus $s$ capiatur infinitus. Curva ergo erit ex spiralium genere, ita ut infinitis præactis spiris in certo quodam puncto tanquam centro convolvatur, quod punctum ex hac constructione invenire difficillimum videtur. Non exiguum ergo analysis incrementum capere existimanda erit, si quis methodum invenerit, cujus ope, saltem vero proxime, valor horum integralium $\int ds\sin.\frac{ss}{2aa}$ \& $\int ds\cos.\frac{ss}{2aa}$ assignari posset, casu quo $s$ ponitur infinitum; quod Problema non indignum videtur, in quo Geometræ vires suas exerceant.

\originalsection{53}\sourcepage{277}Sit $2aa=bb$, \& cum sit
\[\begin{aligned}
\sin.\frac{ss}{bb}&=\frac{s^2}{b^2}-\frac{s^6}{1.2.3b^6}+\frac{s^{10}}{1.2.3.4.5b^{10}}-\frac{s^{14}}{1.2\ldots7b^{14}}+\&c.\\
\cos.\frac{ss}{bb}&=1-\frac{s^4}{1.2b^4}+\frac{s^8}{1.2.3.4b^8}-\frac{s^{12}}{1.2.3.4.5.6b^{12}}+\&c.
\end{aligned}\]
coordinatæ $x$ \& $y$ curvæ quæsitæ commode per series infinitas exprimi poterunt: erit enim
\[\begin{aligned}
x&=\frac{s^3}{1.3b^2}-\frac{s^7}{1.2.3.7b^6}+\frac{s^{11}}{1.2.3.4.5.11b^{10}}-\frac{s^{15}}{1.2\ldots7.15b^{14}}+\&c.\\
y&=s-\frac{s^5}{1.2.5b^4}+\frac{s^9}{1.2.3.4.9b^3}-\frac{s^{13}}{1.2.3\ldots6.13b^{12}}+\&c.
\end{aligned}\]
ex quibus seriebus vehementer convergentibus, nisi arcus $s$ assumatur valde magnus, valores coordinatarum $x$ \& $y$ vero proxime satis expedite determinari possunt. Verum cujusmodi valores $x$ \& $y$ acquirant, si ponatur arcus $s$ infinite magnus, ex his seriebus nullo modo concludi potest.

\originalsection{54}Quoniam igitur positio infiniti loco $s$ facienda maximam parit difficultatem; huic quidem incommodo sequenti modo remedium afferri potest. Ponatur $\frac{ss}{bb}=v$, ut sit $s=b\sqrt v$, erit $ds=\frac{b\,dv}{2\sqrt v}$, fietque, $x=\frac b2\int\frac{dv}{\sqrt v}\sin.v$, \& $y=\frac b2\int\frac{dv}{\sqrt v}\cos.v$. Nunc autem dico valores debitos pro $x$ \& $y$, si ponatur $s=\infty$, inventum iri ex his formulis integralibus,
\[\begin{aligned}
x&=\frac b2\int dv\left(\frac1{\sqrt v}-\frac1{\sqrt{\pi+v}}+\frac1{\sqrt{2\pi+v}}-\frac1{\sqrt{3\pi+v}}+\&c.\right)\sin.v,\\
y&=\frac b2\int dv\left(\frac1{\sqrt v}-\frac1{\sqrt{\pi+v}}+\frac1{\sqrt{2\pi+v}}-\frac1{\sqrt{3\pi+v}}+\&c.\right)\cos.v,
\end{aligned}\]
si post integrationem ponatur $v=\pi$, denotante $\pi$ angulum duobus rectis æqualem. Hoc ergo modo positio infiniti quidem evitatur, contra vero series infinita $\frac1{\sqrt v}-\frac1{\sqrt{\pi+v}}+\frac1{\sqrt{2\pi+v}}-\&c.$ in calculum introducitur, cujus summa cum adhuc lateat, resolutio hujus nodi maximæ adhuc difficultati est obnoxia.

\originalsection{55}\sourcepage{278}Tradita jam methodo investigandi curvaturam cujusque laminæ elasticæ, si ea ab una vi in dato loco applicata sollicitetur; conveniet quoque curvaturam a pluribus, imo infinitis, potentiis laminæ elasticæ inductam indagare. Quoniam vero nondum constat, cujusmodi expressio his casibus futura sit vel maxima vel minima; methodo utar tantum directa, quo ex ipsa solutione fortasse proprietas ea, quæ est maxima vel minima erui queat. \marginnote{\hyperlink{addfig:18}{Fig.~18}}Sit igitur lamina elastica, naturaliter recta, in statum $AmM$ redacta, primum a viribus finitis $P$ \& $Q$ secundum directiones $CE$ \& $CF$ inter se normales sollicitantibus, tum vero a viribus infinite parvis singulis laminæ elementis $m\mu$ applicatis, \& secundum directiones $mp$ \& $mq$ illis $CE$ \& $CF$ parallelas trahentibus; quibus positis requiritur natura curvæ $AmM$ laminæ inductæ.

\originalsection{56}Sumatur recta $FCA$ producta pro axe, ponatur $AC=c$, \& vocetur abscissa $AP=x$, applicata $PM=y$, arcus curvæ $AM=s$, \& radius osculi in $M=R$. Sit elasticitas laminæ absoluta constans $=Ekk$; atque summa momentorum ex omnibus viribus sollicitantibus respectu puncti $M$ ortorum æqualis esse debet $\frac{Ekk}R$. Primum quidem a vi finita $P$ in directione $CE$ trahente oritur momentum $=P(c+x)$, in eam plagam agens qua vis elastica æquilibratur. Momentum autem ex altera vi $Q$ ortum, nempe $Qy$, in contrariam plagam tendit, ex quo ex viribus finitis $P$ \& $Q$ conjunctim oritur momentum $P(c+x)-Qy$. Jam consideretur quodvis elementum laminæ intermedium $m\mu$, cujus respondens abscissa $Ap$ ponatur $=\zeta$, \& applicata $pm=\eta$, sit autem vis elementum $m\mu$ in directione $mp$ urgens $=dp$, \& vis urgens in directione $mq=dq$; erit momentum ex his viribus pro puncto $M$ ortum $=(x-\zeta)dp-(y-\eta)dq$.

\originalsection{57}Ad summam ergo omnium horum momentorum inveniendam, punctum $M$, ac proinde $x$ \& $y$, tantisper pro constantibus haberi debent, dum solæ coordinatæ $\zeta$ \& $\eta$ cum viribus $dp$ \& $dq$ tanquam variabiles spectantur. Erit ergo summa
\sourcepage{279}momentorum a viribus arcum $Am$ sollicitantibus ortorum $=xp-\int\zeta dp-yp+\int\eta dq$; ubi $p$ exprimit summam omnium virium arcum $AM$ in directionibus applicatis $pm$ parallelis sollicitantium, \& $q$ summam omnium virium arcum $Am$ in directionibus axi $Ap$ parallelis sollicitantium. At est $\int\zeta dp=\zeta p-\int p\,d\zeta$ \& $\int\eta dq=\eta q-\int q\,d\eta$; unde fit summa momentorum ex viribus arcui $Am$ applicatis ortorum $=(x-\zeta)p+\int p\,d\zeta-(y-\eta)q-\int q\,d\eta$. Promoveatur jam punctum $m$ in $M$ usque, fietque $\zeta=x$, $\eta=y$, \& $d\zeta=dx$ atque $d\eta=dy$; unde summa omnium momentorum per totum arcum $AM$ sumptorum erit $=\int pdx-\int qdy$. Quocirca obtinebitur pro curva quæsita hæc æquatio $\frac{Ekk}R=P(c+x)-Qy+\int pdx-\int qdy$, ubi ergo $p$ exprimit summam omnium virium verticalium seu in directionibus applicatarum $MP$ agentium, \& $q$ summam omnium virium horizontalium seu in directionibus $MQ$ axi $AP$ parallelis agentium per totum arcum $AM$.

\originalsection{58}Si formulæ $pdx$ \& $qdy$ integrationem non admittant; tum æquatio inventa per differentiationem ab his formulis integralibus liberari debebit, unde habebitur ista æquatio:
\[\frac{-EkkdR}{RR}=Pdx-Qdy+pdx-qdy.\]
Sin autem nec $p$ nec $q$ per expressiones finitas exhiberi possint; quippe quæ jam exprimunt summas infinitarum virium infinite parvarum, tum per ulteriorem differentiationem valores finiti $p$ \& $q$ exterminari debebunt, ut tantum insint $dp$ \& $dq$ cum differentio-differentialibus $ddp$ \& $ddq$. Orietur autem, post primam differentiationem, $-Ekkd.\frac{dR}{RRdx}=dp-(Q+q)\times d.\frac{dy}{dx}-\frac{dy}{dx}dq$. Sit $\frac{dy}{dx}=\omega$, eritque denuo æquatione differentiata:
\[-Ekkd.\frac{d.\frac{dR}{RRdx}}{d\omega}=d.\frac{dp}{d\omega}-2dq-\omega d.\frac{dq}{d\omega},\]
quæ æquatio ad differentialia quarti ordinis ascendit.

\originalsection{59}\sourcepage{280}Sint laminæ, loco potentiarum verticalium \& horizontalium $p$ \& $q$, in singulis punctis $M$ applicatæ duæ potentiæ altera normalis $MN=dv$ \& altera tangentialis $MT=dt$. Hinc erit $dp=\frac{dx\,dv}{ds}+\frac{dy\,dt}{ds}$ \& $dq=\frac{dx\,dt}{ds}-\frac{dy\,dv}{ds}$; \&, ob $dy=\omega dx$ \& $ds=dx\sqrt{1+\omega\omega}$, habebitur $dp=\frac{dv}{\sqrt{1+\omega\omega}}+\frac{\omega dt}{\sqrt{1+\omega\omega}}$, \& $dq=\frac{dt}{\sqrt{1+\omega\omega}}-\frac{\omega dt}{\sqrt{1+\omega\omega}}$: quibus in præced. §. æquatione ultima substitutis, proveniet sequens æquatio,
\[-Ekkd.\frac{d.\frac{dR}{RRdx}}{dv}=\frac{-dt}{\sqrt{1+\omega\omega}}+\frac{2\omega dv}{\sqrt{1+\omega\omega}}+\sqrt{1+\omega\omega}d.\frac{dv}{d\omega},\]
quæ multiplicata per $\sqrt{1+\omega\omega}$ fit integrabilis: posito enim, brevitatis gratia, $z=\frac{dR}{RRdx}$, reperietur integrale,
\[\begin{aligned}
A-t+\frac{dv(1+\omega\omega)}{d\omega}&=-Ekk\left(\frac{dz\sqrt{1+\omega\omega}}{d\omega}-\frac{\omega z}{\sqrt{1+\omega\omega}}+\frac1{2RR}\right)\\
&=-Ekk\left(\frac{1+\omega\omega}{d\omega}d.\frac{dR}{RRdx\sqrt{1+\omega\omega}}+\frac1{2RR}\right).
\end{aligned}\]
Cum vero sit $R=\frac{-(1+\omega\omega)^{3:2}dx}{d\omega}$, erit $d\omega=\frac{-(1+\omega\omega)^{3:2}dx}R$; quo loco $d\omega$ valore substituto, habebitur:
\[A-t-\frac{Rdv}{ds}=-Ekk\left(\frac1{2RR}-\frac R{ds}d.\frac{dR}{RRds}\right),\]
ob $dx\sqrt{1+\omega\omega}=ds$. Quocirca æquatione ordinata, pro curva quæsita orietur hæc æquatio
\[t+\frac{Rdv}{ds}-A=Ekk\left(\frac1{2RR}-\frac R{ds}d.\frac{dR}{RRds}\right).\]

\originalsection{60}Primum quidem manifestum est, si vis elastica $Ekk$ evanescat, laminam transmutari in filum perfecte flexile; atque hinc in his æquationibus continentur omnes curvæ, quas filum perfecte flexile a viribus quibuscunque sollicitatum formare potest. Sic si filum a propria gravitate tantum deorsum sollicitatur, erit
\sourcepage{281}$q=0$, \& $p$ exprimet pondus funis $AM$, eritque ergo $\frac{pdx}{dy}=Q=\text{constanti}$, facto $P=0$, quæ est æquatio generalis pro omnis generis Catenariis. Sin autem filum perfecte flexile, in singulis punctis, a viribus quarum directiones sunt normales ad ipsam curvam sollicitetur, ita ut in puncto $M$ filum sollicitetur secundum directionem $MN$, vi $=dv$; ob $t=0$, erit $\frac{Rdv}{ds}=A=\text{constanti}$, quæ est proprietas generalis curvarum Velariarum, Linteariarum, omniumque in quibus hujusmodi sollicitationes locum habent.

\originalsection{61}Ad laminas elasticas autem revertor, de quibus mox ista quæstio præ ceteris notatu digna se offert, cujusmodi figuram accipiat lamina elastica proprio pondere incurvata. Sit $AmM$ hæc curva quæ quæritur, \& quia solæ vires verticales a gravitate ortæ urgent, fiet $P=0$, $Q=0$, $q=0$, \& $p$ exprimet pondus laminæ $AM$. Quare si $F$ sit pondus laminæ longitudinis $a$; quia lamina uniformis assumitur, erit $p=\frac{Fs}a$; unde curvæ natura hac exprimetur æquatione $\frac{-EkkdR}{RR}=\frac{Fsdx}a$. Sit amplitudo curvæ $\int\frac{ds}R=u$, erit $R=\frac{ds}{du}$, \& $dx=ds\sin.u$; unde, sumpto elemento $ds$ constante, reperietur æquatio $sds\sin.u+\frac{Eakk}F.\frac{ddu}{ds}=0$, quæ autem, quantum primo intuitu patet, ulterius reduci nequit.

\originalsection{62}In primis autem notari meretur curva, quam fluidum altitudinis quasi infinitæ laminæ elasticæ inducit. \marginnote{\hyperlink{addfig:19}{Fig.~19}}Sit $AMB$ figura hæc quæ quæritur, \& posito $AP=x$, $PM=y$, $AM=s$; elementum $Mm$ in directione normali $MN$ urgebitur vi ipsi $ds$ proportionali; unde erit $dv=nds$, \& $dt=0$. Hinc oritur vis verticalis $dp=ndx$, \& horizontalis $dq=-ndy$; ex quibus statim fit $p=nx$ \& $q=-ny$; ideoque in æquatione prima fiet $\frac{Ekk}R=P(c+x)-Qy+\frac12nxx+\frac12nyy$.
\sourcepage{282}Coordinatæ vero $x$ \& $y$ ita quantitatibus constantibus augeri diminuive possunt, ut æquatio pro curva hujusmodi faciem acquirat $xx+yy=A+\frac BR$. Hæc autem æquatio si multiplicetur per $xdx+ydy$, fiet integrabilis; est enim
\[\int\frac{xdx+ydy}R=-\int\frac{x+y\omega}{(1+\omega\omega)^{3:2}}d\omega,\quad[\text{posito }dy=\omega dx]=\frac{y-\omega x}{\sqrt{1+\omega\omega}}=\frac{ydx-xdy}{ds}.\]
Hanc ob rem, post integrationem constantibus mutatis, prodibit $(xx+yy)^2=A(xx+yy)+\frac{B(ydx-xdy)}{ds}+C$. Sit $\sqrt{xx+yy}=z$ \& $y=uz$; erit $x=z\sqrt{1-uu}$; unde $ydx-xdy=\frac{-zzdu}{\sqrt{1-uu}}$, \& $ds=\sqrt{dz^2+\frac{zzdu^2}{1-uu}}$. Ergo posito $\frac{du}{\sqrt{1-uu}}=dr$, erit $z^4-Az^2-C=\frac{-Bzzdr}{\sqrt{dz^2+zzdr^2}}$; hincque
\[dr=\frac{du}{\sqrt{1-uu}}=\frac{dz(z^4-Az^2-C)}{z\sqrt{B^2zz-(z^4-Az^2-C)^2}}.\]
Curva hæc ergo, si fuerit $A=0$ \& $C=0$, erit algebraica; habebitur enim hæc æquatio $\frac{du}{\sqrt{1-uu}}=\frac{zzdz}{\sqrt{B^2-z^6}}=\frac{3zzdz}{3\sqrt{a^6-z^6}}$, quæ integrata dat $\operatorname{Asin}.u=\frac13\operatorname{Asin}.\frac{z^3}{a^3}$, seu $\frac{z^3}{a^3}=3u-4u^3=\frac{3y}z-\frac{4y^3}{z^3}$; unde hæc resultat æquatio $z^6=3a^3yzz-4a^3y^3$; seu, ob $zz=xx+yy$, hæc $x^6+3x^4y^2+3xxy^4+y^6=3a^3xxy-a^3y^3$.

\originalsection{63}Ex his etiam motus oscillatorius laminarum elasticarum utcunque ad motum comparatarum definiri potest: quod argumentum profecto dignissimum primum excolere cœpit Vir Celeb. \textit{Daniel Bernoulli}, mihique, jam ante complures annos, Problema de oscillationibus laminæ elasticæ altero termino parieti firmo infixæ determinandis proposuit, cujus Solutionem
\sourcepage{283}exhibui in \textit{Comment. Petropol.} Tomo VII. Ex hoc autem tempore, cum mihi commodius hoc Problema tractare contigit; tum etiam per commercium cum Celeb. Bernoullio plures accesserunt aliæ quæstiones \& considerationes, quarum enodationem, ob materiæ affinitatem, hic adjungam. Quando autem motus vibratorius est satis promtus, tum simul a lamina vibrante sonus editur, cujus tenor ac relatio ad alios, ope doctrinæ de sonis, ex his principiis determinabitur. Et quoniam sonorum indoles facillime ad experimenta revocatur; hoc ipso consensus calculi cum veritate explorari, atque adeo Theoria confirmari poterit: quo pacto, cognitio nostra circa naturam corporum elasticorum non parum amplificabitur.

\originalsection{64}Primum autem monendum est, hic tantum circa oscillationes minimas quæstionem institui; atque adeo intervallum, per quod lamina inter oscillandum excurrit, esse quasi infinite parvum. Neque vero, hac limitatione, usus \& applicatio quicquam diminuitur: non solum enim oscillationes, si per majora spatia fierent, isochronismo destituerentur; sed etiam sonorum distinctorum formatio, ad quam hic potissimum spectamus, minimas oscillationes requirit. Considero igitur hic primum laminam elasticam uniformem naturaliter rectam, cujus alter terminus $B$ pavimento immobili firmiter sit infixus, ita ut lamina sibi relicta situm teneat rectum $BA$. \marginnote{\hyperlink{addfig:20}{Fig.~20}}Sit hujus laminæ longitudo $AB=a$, ejusque elasticitas absoluta in singulis locis $=Ekk$; ab ejus vero pondere vel mentem revocamus, vel infixionem ejusmodi statuimus, ut ejus status a gravitate turbari nequeat.

\originalsection{65}Jam lamina hæc a vi quacunque impulsa vibrationes peragat minimas, circa statum naturalem $BA$ utrinque excurrendo per minima intervalla $Aa$. Sitque $BMa$ status quispiam, quem lamina inter oscillandum tenet; qui quoniam infinite parum tantum distat a statu naturali $BPA$, rectæ $MP$, $Aa$ simul repræsentabunt vias, quas laminæ puncta $M$ \& $a$ percurrunt, vel potius hæ rectæ ad vias veras rationem habebunt a ratione æqualitatis infinite parum discrepantem. Ad motum autem oscillatorium determinandum, absolute necesse est naturam \sourcepage{284}curvæ $BMa$, quam lamina inter oscillandum induit, nosse. Sit igitur $AP=x$, $PM=y$, arcus $aM=s$, \& radius osculi in $M=R$; \& intervallum minimum $Aa=b$; atque, ex conditione memorata, erit arcus $s$ proxime æqualis abscissæ $x$, ac proinde pro $ds$ sumi poterit $dx$: præ $dx$ enim evanescet $dy$. Et cum, posito $dx$ constante, sit generatim radius osculi $=\frac{ds^3}{dxddy}$, erit præsenti casu $R=\frac{dx^2}{ddy}$; nam curva $BMa$ convexitatem axi $BA$ obvertit, \& quia lamina in $B$ muro firmiter est infixa, erit recta $AB$ tangens curvæ in puncto $B$.

\originalsection{66}His positis, tam ad naturam curvæ $BMa$ quam ad ipsum motum oscillatorium determinandum, sit $f$ longitudo penduli simplicis isochroni; oscillationes enim minimas esse isochronas, cum natura rei declarat, tum ipse calculus instituendus monstrabit. Acceleratio ergo, qua laminæ punctum $M$ versus $P$ urgetur, erit $=\frac{PM}f=\frac yf$. Quare si massa totius laminæ ponatur $=M$, quæ per ejus pondus exprimitur; erit elementi $Mm=ds=dx$ massa $=\frac{Mdx}a$; unde vis motrix elementum $Mm$ in directione $MP$ sollicitans erit $=\frac{Mydx}{af}$; sicque vires, quibus singulæ laminæ particulæ ad motum actu cientur, innotescunt, cum ex ipsa curva $BMa$, tum ex longitudine penduli simplicis isochroni $f$. Quoniam vero lamina a vi elastica revera ad motum incitatur; ex hac cognita vicissim \& natura curvæ $BMa$, \& longitudo penduli simplicis isochroni determinabitur.

\originalsection{67}Quoniam ergo lamina perinde movetur, ac si singulis ipsius elementis $Mm$ in directione $MP$ vires essent applicatæ $=\frac{Mydx}{af}$; sequitur, si laminæ singulis elementis $Mm$ in directionibus contrariis $M\pi$ æquales vires $\frac{Mydx}{af}$ applicarentur, laminam in statu $BMa$ æquilibrari. Hinc lamina inter oscillandum eandem curvaturam subibit, quam indueret quieta, si in
\sourcepage{285}singulis punctis $M$ sollicitaretur viribus $\frac{Mydx}{af}$ in directionibus $M\pi$. Per regulam ergo supra §.~56 inventam, colligantur omnes hæ vires per arcum $aM$ applicatæ, atque prodibit summa $=\frac M{af}\int ydx$, quæ ibi in locum ipsius $p$ substitui debet. Quare cum vires reliquæ $P$, $Q$, \& $q$, quæ ibi habebantur, evanescant, natura curvæ exprimetur æquatione $\frac{Ekk}R=\int pdx$: unde habebitur $\frac{Ekk}R=\frac M{af}\int dx\int ydx$. Cum vero sit $R=\frac{dx^2}{ddy}$, erit $\frac{Ekkddy}{dx^2}=\frac M{af}\int dx\int ydx$; \& differentiando $\frac{Ekkd^3y}{dx^2}=\frac{Mdx}{af}\times\int ydx$: denuoque differentiando prodibit ista æquatio differentialis quarti ordinis. $Ekkd^4y=\frac{Mydx^4}{af}$.

\originalsection{68}Hac ergo æquatione \& natura curvæ $BMa$ exprimitur, \& ex eadem, si ad casum oblatum accommodetur, longitudo $f$ determinabitur; qua cognita, ipse motus oscillatorius innotescet. Ante omnia autem hanc æquationem integrari oportet: quæ cum pertineat ad id æquationum differentialium altiorum graduum genus, cujus integrationem generalem exhibui in \textit{Miscell. Berol.} Volumine VII, hinc sequens æquatio integralis reperietur; ponendo brevitatis ergo $\frac{Ekkaf}M=c^4$; prodibit scilicet
\[y=Ae^{\frac xc}+Be^{-\frac xc}+C\sin.\frac xc+D\cos.\frac xc,\]
ubi $e$ denotat numerum cujus logarithmus hyperbolicus est $=1$; \& $\sin.\frac xc$ \& $\cos.\frac xc$ denotant sinum \& cosinum arcus $=\frac xc$ in circulo, cujus radius $=1$, assumpti. Tum vero $A$, $B$, $C$, \& $D$ sunt quatuor constantes arbitrariæ per quadruplicem integrationem introductæ, quas ex accommodatione calculi ad præsentem casum definire oportet.

\originalsection{69}Determinatio autem constantium sequenti modo instituitur. \sourcepage{286}Primum posito $x=0$, fieri debet $y=b$; hinc ergo oritur ista æquatio, $b=A+B+D$, quæ est \textit{prima}.

Secundo, cum sit $\frac{c^4ddy}{dx^2}=\int dx\int ydx$; facto $x=0$, fieri debet $\frac{ddy}{dx^2}=0$; at est $\frac{ddy}{dx^2}=\frac A{cc}e^{\frac xc}+\frac B{cc}e^{-\frac xc}-\frac C{cc}\sin.\frac xc-\frac D{cc}\cos.\frac xc$; unde oritur hæc æquatio \textit{secunda}, $0=A+B-D$.

Tertio, cum sit $\frac{c^4d^3y}{dx^3}=\int ydx$; posito $x=0$, simul $\frac{d^3y}{dx^3}$ evanescere debet: quia ergo erit $\frac{c^3d^3y}{dx^3}=Ae^{\frac xc}-Be^{-\frac xc}-C\cos.\frac xc+D\sin.\frac xc$: prodit æquatio \textit{tertia}, $0=A-B-C$.

Quarto autem, si ponatur $x=a$, applicata $y$ evanescit, unde obtinebitur æquatio \textit{quarta}, $0=Ae^{\frac ac}+Be^{-\frac ac}+C\sin.\frac ac+D\cos.\frac ac$.

Quinto, quia $AB$ est tangens curvæ in puncto $B$; facto $x=a$, fieri debet $\frac{dy}{dx}=0$: unde prodit æquatio \textit{quinta},
\[0=Ae^{\frac ac}-Be^{-\frac ac}+C\cos.\frac ac-D\sin.\frac ac.\]
Ex his ergo quinque æquationibus, primum quatuor constantes $A$, $B$, $C$, $D$ definientur; tum vero, in quo cardo rei versatur, determinabitur valor ipsius $c=\sqrt[4]{\frac{Ekkaf}M}$; ex quo longitudo penduli simplicis isochroni $f$ elicietur; quo ipso, durationes oscillationum cognoscentur.

\originalsection{70}Ex æquationibus secunda \& tertia, constantes $C$ \& $D$ ex $A$ \& $B$ ita definientur, ut sit $C=A-B$, \& $D=A+B$.
\sourcepage{287}qui valores in æquationibus quarta \& quinta substituti dabunt
\[\begin{aligned}
0&=Ae^{\frac ac}+Be^{-\frac ac}+(A-B)\sin.\frac ac+(A+B)\cos.\frac ac,\\
0&=Ae^{\frac ac}-Be^{-\frac ac}+(A-B)\cos.\frac ac-(A+B)\sin.\frac ac;
\end{aligned}\]
ex quibus eruitur,
\[\frac AB=\frac{-e^{-\frac ac}+\sin.\frac ac-\cos.\frac ac}{e^{\frac ac}+\sin.\frac ac+\cos.\frac ac}=\frac{e^{-\frac ac}+\cos.\frac ac+\sin.\frac ac}{e^{\frac ac}+\cos.\frac ac-\sin.\frac ac},\]
unde obtinebitur hæc æquatio,
\[0=2+(e^{\frac ac}+e^{-\frac ac})\cos.\frac ac,\]
seu $e^{\frac{2a}c}\cos.\frac ac+2e^{\frac ac}+\cos.\frac ac=0$, quæ dat $e^{\frac ac}=\frac{-1\pm\sin.\frac ac}{\cos.\frac ac}$. Cum igitur $e^{\frac ac}$ sit quantitas affirmativa, cosinus anguli $\frac ac$ erit negativus; ideoque angulus $\frac ac$ recto major.

\originalsection{71}Ex hac æquatione intelligitur dari infinitos angulos $\frac ac$ quæsito satisfacientes, ex quibus infiniti diversi modi oscillationum ejusdem laminæ oriuntur. Curva enim in uno pluribusve punctis axem $AB$ secare potest, antequam in $B$ axem tangat; ex quo ejusdem laminæ plures, imo infiniti, oscillandi modi æque sunt possibiles. Cum igitur hic inprimis contemplemur casum, quo $B$ primum est punctum, ubi lamina ab axe $AB$ tangitur; huic casui satisfaciet minimus angulus $\frac ac$ æquationem \sourcepage{288}inventam resolvens; qui angulus cum sit recto major, ponatur $\frac ac=\frac12\pi+\varphi$; existente $\varphi$ angulo recto minore. Hinc, ob $\sin.\frac ac=\cos.\varphi$, \& $\cos.\frac ac=-\sin.\varphi$, obtinebitur duplex æquatio,
\[e^{\frac ac}=\frac{1\mp\cos.\varphi}{\sin.\varphi}.\]
quæ præbet, vel $e^{\frac ac}=\operatorname{tang}.\frac12\varphi$, vel $e^{\frac ac}=\cot.\frac12\varphi$, quarum posterior minorem dabit valorem pro angulo $\varphi$; quæ ergo ad casum propositum erit accommodata.

\originalsection{72}Sequentes possibiles oscillationum modi reperientur, si pro $\frac ac$ ponantur anguli duobus rectis majores, tribus vero minores. Sic posito $\frac ac=\frac32\pi-\varphi$, erit $\sin.\frac ac=-\cos.\varphi$, \& $\cos.\frac ac=-\sin.\varphi$; unde fit $e^{\frac ac}=\frac{1\pm\cos.\varphi}{\sin.\varphi}$, seu, vel $e^{\frac ac}=\operatorname{tang}.\frac12\varphi$, vel $e^{\frac ac}=\cot.\frac12\varphi$. Simili modo alii oscillationum modi reperientur, ponendo $\frac ac=\frac52\pi+\varphi$; $\frac ac=\frac72\pi-\varphi$, \&c. Ex quibus omnibus, si sumantur logarithmi hyperbolici, orientur sequentes æquationes:
\[\begin{array}{rl@{\qquad}rl}
\mathrm{I}.&\frac12\pi+\varphi=l\cot.\frac12\varphi;&\mathrm{II}.&\frac12\pi+\varphi=l\operatorname{tang}.\frac12\varphi;\\
\mathrm{III}.&\frac32\pi-\varphi=l\cot.\frac12\varphi;&\mathrm{IV}.&\frac32\pi-\varphi=l\operatorname{tang}.\frac12\varphi;\\
\mathrm{V}.&\frac52\pi+\varphi=l\cot.\frac12\varphi;&\mathrm{VI}.&\frac52\pi+\varphi=l\operatorname{tang}.\frac12\varphi;\\
\mathrm{VII}.&\frac72\pi-\varphi=l\cot.\frac12\varphi;&\mathrm{VIII}.&\frac72\pi-\varphi=l\operatorname{tang}.\frac12\varphi;\\
&&\&c.
\end{array}\]
Harum autem æquationum tertia congruit cum secunda; posito enim $\frac12\varphi=\frac12\pi-\frac12\theta$, ut sit $\cot.\frac12\varphi=\operatorname{tang}.\frac12\theta$; tertia transit in $\frac12\pi+\theta=l\operatorname{tang}.\frac12\theta$, quæ est ipsa secunda. Simili modo quarta congruit cum prima: tum quinta \& octava inter se
\sourcepage{289}congruunt; atque sexta cum septima. Quamobrem sequentes tantum prodibunt æquationes diversæ:
\[\begin{array}{rl}
\mathrm{I}.&\frac12\pi+\varphi=l\cot.\frac12\varphi\\
\mathrm{II}.&\frac12\pi+\varphi=l\operatorname{tang}.\frac12\varphi\\
\mathrm{III}.&\frac52\pi+\varphi=l\cot.\frac12\varphi\\
\mathrm{IV}.&\frac52\pi+\varphi=l\operatorname{tang}.\frac12\varphi\\
\mathrm{V}.&\frac92\pi+\varphi=l\cot.\frac12\varphi\\
\mathrm{VI}.&\frac92\pi+\varphi=l\operatorname{tang}.\frac12\varphi\\
&\&c.
\end{array}\]

\originalsection{73}\footnote{The source prints 72 here; the consecutive section number 73 is used in this edition.}Logarithmus autem hyperbolicus tangentis vel cotangentis cujuspiam anguli reperitur, sumendo logarithmum Tabularem, indeque auferendo logarithmum sinus totius, atque residuum multiplicando per $2{,}302585092994$; qui labor ut sublevetur denuo logarithmis uti conveniet. Sit $u$ logarithmus hyperbolicus tangentis seu cotangentis anguli $\frac12\varphi$, qui quæritur; sumatur ex Tabulis logarithmus ejusdem tangentis cotangentisve, qui logarithmo sinus totius multatus ponatur $=v$. Cum ergo sit $u=2{,}302585092994\times v$; erit, sumendis logarithmis vulgaribus,
\[lu=lv+0{,}3622156886.\]
Hoc logarithmo invento, cum sit $u=\frac n2\pi+\varphi$, erit $lu=l(\frac n2\pi+\varphi)$. Ad hoc evolvendum, angulus $\varphi$ in partibus radii exprimi debet, quemadmodum \& $\pi$ eodem modo exprimitur, dum est $\pi=3{,}1415926535$, ac propterea $\frac12\pi=1{,}57079632679$. Angulus autem $\varphi$ eodem modo exprimitur, si is in minuta secunda convertatur, atque ab hujus numeri logarithmo subtrahatur constanter $5{,}3144251332$; sic enim prodibit $l\varphi$, ex quo ad numeros regrediendo valor ipsius $\varphi$ eruitur. Erit autem constanter pro unoquoque oscillationum genere $\frac ac=u=\frac n2\pi+\varphi$.

\originalsection{74}His circa calculum instituendum monitis, per approximationes \sourcepage{290}valor anguli $\varphi$ pro quovis oscillationum genere non difficulter eruetur. Tribuendo enim pro lubitu ipsi $\varphi$ valores aliquot, \& per calculum determinando, \& $\frac n2\pi+\varphi$, \& $l\begin{smallmatrix}\operatorname{tang}.\\\cot.\end{smallmatrix}\frac12\varphi$, mox valor ipsius $\varphi$ prope verus cognoscetur. Quod si autem habeantur limites anguli $\varphi$ utcunque remoti, statim invenientur limites propiores, ex hisque tandem verus valor ipsius $\varphi$. Sic pro æquatione prima $\frac ac=\frac12\pi+\varphi=l\cot.\frac12\varphi$, sequentes limites anguli $\varphi$ erui $17^\circ$, $26'$, \& $17^\circ$, $27'$, ex quibus per sequentem calculum verus valor ipsius $\varphi$ obtinebitur.
\[\renewcommand{\arraystretch}{1.15}\begin{array}{r@{\;=\;}r|r}
\varphi&17^\circ,26',0''&17^\circ,27',0''\\
\text{in min. sec.}&62760''&62820''\\
\log.&4{,}7976829349&4{,}7980979321\\
\text{subtr.}&5{,}3144251332&5{,}3144251332\\\hline
l\varphi&9{,}4832578017&9{,}4836727989\\
\varphi&0{,}3042690662&0{,}3045599545\\
\frac12\pi&1{,}5707963268&1{,}5707963268\\\hline
\frac12\pi+\varphi&1{,}8750653930&1{,}8753562813\\\hline
\frac12\varphi&8^\circ,43',0''&8^\circ,43',30''\\
l\cot.\frac12\varphi&10{,}8144034109&10{,}8139819342\\
v&0{,}8144034109&0{,}8139819342\\
lv&9{,}9108395839&9{,}9106147660\\
\text{add.}&0{,}3622156886&0{,}3622156886\\\hline
lu&0{,}2730552725&0{,}2728304546\\
u&1{,}8752331540&1{,}8742626675\\\hline
\text{diff.}&+1677610&-10936138
\end{array}\]
Ex his ergo utriusque limitis erroribus concluditur fore $\varphi=17^\circ$, $26'$, $7''\frac{98}{100}$, \& $\frac12\pi+\varphi$, seu $\frac ac=107^\circ$, $26'$, $7''\frac{98}{100}$. Cum vero in minutis secundis sit $\varphi=62767{,}98$, erit
\sourcepage{291}
\[\begin{array}{r@{\;=\;}r}
l\varphi&4{,}7977381525\\
\text{subtr.}&5{,}3144251332\\\cline{2-2}
\multicolumn{1}{r}{}&9{,}483130193\\
\text{ergo }\varphi&0{,}3043077545\\
\text{add. }\frac12\pi&1{,}5707963268\\\cline{2-2}
\frac ac&1{,}8751040813
\end{array}\]
quo invento, erit $\frac AB=\operatorname{tang}.\frac12\varphi=0{,}1533390624$. Reperitur ergo ratio constantium $A$ \& $B$, ex quibus \& ratio reliquarum constantium $C$ \& $D$ ad illas cognoscetur.

\originalsection{75}Restat adhuc prima æquatio $b=A+B+D$, quæ, ob $D=A+B$, abit in $b=2A+2B$; ideoque $A+B=\frac12b$: cum ergo sit $\frac AB=\operatorname{tang}.\frac12\varphi$, fiet $B(1+\operatorname{tang}.\frac12\varphi)=\frac12b$, \& $B=\frac b{2+2\operatorname{tang}.\frac12\varphi}$. Unde ex $\operatorname{tang}.\frac12\varphi=0{,}1533390624$ singulæ æquationis constantes sequenti modo determinabuntur:
\[\begin{aligned}
\frac Ab&=\frac{\operatorname{tang}.\frac12\varphi}{2(1+\operatorname{tang}.\frac12\varphi)}=\frac{0{,}1533390624}{2{,}3066781248},\\
\frac Bb&=\frac1{2(1+\operatorname{tang}.\frac12\varphi)}=\frac{1{,}0000000000}{2{,}3066781248},\\
\frac Cb&=\frac{-1+\operatorname{tang}.\frac12\varphi}{2(1+\operatorname{tang}.\frac12\varphi)}=\frac{-0{,}8466609376}{2{,}3066781248},\\
\frac Db&=\frac{1+\operatorname{tang}.\frac12\varphi}{2(1+\operatorname{tang}.\frac12\varphi)}=\frac{1{,}1533390624}{2{,}3066781248},
\end{aligned}\]
quibus inventis, natura curvæ $aMB$, quam lamina inter oscillandum induit, hac exprimetur æquatione
\[\frac yb=\frac Abe^{\frac xc}+\frac Bbe^{-\frac xc}+\frac Cb\sin.\frac xc+\frac Db\cos.\frac xc.\]

\originalsection{76}Quod autem ad oscillationum velocitatem attinet, ea \sourcepage{292}ex æquatione $\frac ac=1{,}8751040813$ cognoscetur. Ponatur brevitatis gratia: $n=1{,}8751040813$ ut sit $a=nc$. Et cum sit $c^4=\frac{Ekkaf}M$, ubi $\frac Ma$ exprimit gravitatem specificam laminæ, \& $Ekk$ elasticitatem absolutam; eo modo, quo hactenus sum usus, erit $a^4=n^4.Ekk.\frac aMf$, ideoque $f=\frac{a^4}{n^4}.\frac1{Ekk}.\frac Ma$, ex quo longitudo penduli simplicis isochroni tenebit rationem compositam ex quadruplicata longitudinis laminæ, simplici gravitatis specificæ, \& inversa elasticitatis absolutæ. Sit $g$ longitudo penduli simplicis singulis minutis secundis oscillantis, ita ut sit $g=3{,}16625$ ped. Rhenani; quia durationes oscillationum sunt in subduplicata ratione pendulorum, tempus unius oscillationis a lamina nostra elastica factæ, erit $=\frac{\sqrt f}{\sqrt g}\text{ secund.}=\frac{aa}{nn}\sqrt{\frac1g.\frac1{Ekk}.\frac Ma}$; unde numerus oscillationum uno minuto secundo editarum erit $=\frac{nn}{aa}\sqrt{g.Ekk.\frac aM}$, qui numerus exprimit soni quem lamina excitat tenorem. Soni ergo a diversis laminis elasticis uno termino muro infixis editi erunt in ratione composita subduplicata elasticitatum absolutarum directe, inversa subduplicata gravitatum specificarum, \& inversa duplicata longitudinum. Quare si duæ laminæ elasticæ tantum longitudine differant, erunt soni reciproce ut quadra longitudinum; scilicet lamina duplo longior edet sonum duabus octavis graviorem. Corda autem tensa duplo longior sonum una tantum octava graviorem edit, si tensio maneat eadem. Ex quo patet sonos laminarum elasticarum longe aliam sequi rationem, atque sonos cordarum tensarum.

\originalsection{77}Quod ad naturam curvæ $aMB$ ultra terminos $a$ \& $B$ continuatæ attinet, primum quidem patet curvam ultra $a$ divergendo ab axe $BA$ continuato progredi. Posito enim $x$ negativo fiet
\sourcepage{293}
\[y=Be^{\frac xc}+Ae^{-\frac xc}-C\sin.\frac xc+D\cos.\frac xc.\]
Hic jam omnes termini sunt affirmativi, quia solus coefficiens $C$ ante obtinuerat valorem negativum; unde dum crescit $x$, etiam $y$ crescere debet, quia numerus $B$ major est quam $A$; atque adeo terminus $Be^{\frac xc}$ prævalet. Quam primum autem $\frac xc$ valorem saltem mediocrem est adeptum; tum iste terminus $Be^{\frac xc}$ tantopere crescit, ut reliqui termini præ eo quasi evanescant. Ob eandem rationem, quia curvæ in $B$ radius osculi non est $=\infty$; est enim $\frac{Ekk}R=\frac M{af}\int dx\int ydx$; curva in $B$ non habebit punctum flexus contrarii, ideoque ad eandem axis $AB$ partem ulterius progredietur; aucta autem abscissa $x$ ultra $AB=a$, tum primus terminus $Ae^{\frac xc}$ mox tam fit magnus, ut reliqui præ eo pro nihilo reputari queant.

\originalsection{78}Hic igitur est primus oscillationum modus inter illos innumerabiles, ad quos eadem lamina se componere potest. \marginnote{\hyperlink{addfig:21}{Fig.~21}}Secundus modus in figura repræsentatus quo lamina in $B$ fixa axem $AB$ in uno puncto $O$ trajicit, deducetur ex æquatione $\frac ac=\frac12\pi+\varphi=l\operatorname{tang}.\frac12\varphi$, seu hac $\frac32\pi-\varphi=l\cot.\frac12\varphi=\frac ac$. Hic per nonnulla tentamina inveni angulum $\varphi$ contineri intra hos limites, $1^\circ$, $2'$, $40''$ \& $1^\circ$, $3'$, $0''$, ex quibus ut ante verus valor ipsius $\varphi$ eruetur.
\sourcepage{294}
\[\renewcommand{\arraystretch}{1.13}\begin{array}{r@{\;=\;}r|r}
\varphi&1^\circ,2',40''&1^\circ,3',0''\\
\text{in min. sec.}&3760''&3780''\\
\log.&3{,}5751878450&3{,}5774917998\\
\text{subtr.}&5{,}3144251332&5{,}3144251332\\\hline
l\varphi&8{,}2607627118&8{,}2630665666\\
\varphi&0{,}0182289944&0{,}0183259571\\
\frac32\pi&4{,}7123889804&4{,}7123889804\\\hline
\frac ac&4{,}6941599860&4{,}6940630233\\
\frac12\varphi&31',20''&31',30''\\
l\cot.\frac12\varphi&2{,}0402552577&2{,}0379511745\\
lv&0{,}3096845055&0{,}3091937748\\
\text{add.}&0{,}3622156886&0{,}3622156886\\\hline
lu&0{,}6719001941&0{,}6714094634\\
u&4{,}6978613391&4{,}6925559924\\
\frac ac&4{,}6941599860&4{,}6940630233\\\hline
\text{Error}&+37013531&-15070309
\end{array}\]
Ex his erroribus concluditur verus valor anguli $\varphi=1^\circ$, $2'$, $54''\frac{213}{1000}$, \& $\frac ac=268^\circ$, $57'$, $5''\frac{787}{1000}$. Cum igitur sit $\varphi=3774{,}213''$ erit
\[\begin{array}{r@{\;=\;}r}
l\varphi&3{,}5768264061\\
\text{subtr.}&5{,}3144251332\\\cline{2-2}
\multicolumn{1}{r}{}&8{,}2624012729\\
\varphi&0{,}0182979009\\
\text{a }\frac32\pi&4{,}7123889804\\\cline{2-2}
\frac ac&4{,}6940910795
\end{array}\]
Sonus ergo laminæ priori modo oscillantis erit ad sonum ejusdem laminæ hoc modo vibrantis, uti est quadratum numeri $1{,}8751040813$ ad quadratum numeri $4{,}6940910795$, hoc est
\sourcepage{295}ut $1$ ad $6{,}266891$, seu in minimis numeris ut $4$ ad $25$, seu ut $1$ ad $6\frac4{15}$. Unde sonus posterior erit ad priorem duplex octava cum quinta \& cum hemitonio fere.

\originalsection{79}Pro sequentibus oscillationum modis ejusdem laminæ elasticæ, quibus lamina inter oscillandum axem $AB$ in duobus pluribusve punctis intersecat, fit angulus $\varphi$ multo minor. Sic pro tertio modo habetur hæc æquatio $\frac52\pi+\varphi=l\cot.\frac12\varphi=\frac ac$. Cum ergo sit $e^{\frac52\pi+\varphi}=\cot.\frac12\varphi$, ob $\varphi$ angulum vehementer parvum, erit $e^{\frac52\pi+\varphi}=e^{\frac52\pi}(1+\varphi+\frac12\varphi^2+\frac16\varphi^3+\&c.)$ \& $\cot.\frac12\varphi=\frac{1-\frac18\varphi\varphi}{\frac12\varphi-\frac1{48}\varphi^3}=\frac2\varphi-\frac\varphi6$. Hinc erit proxime $e^{\frac52\pi}=\frac2\varphi$; ideoque $\varphi=2e^{-\frac52\pi}$, \& propius $\varphi=\frac1{1+\frac12e^{\frac52\pi}}$; unde erit $\frac ac=\frac52\pi+\frac2{e^{\frac12\pi}+2}$; qui posterior terminus est quam-minimus. Simili modo pro quarto oscillationum modo, erit proxime $\frac ac=\frac72\pi-2e^{-\frac72\pi}$, \& ita porro: ob hos alteros terminos evanescentes, ipsius $\frac ac$ valores erunt $\frac92\pi$, $\frac{11}2\pi$, \&c. qui eo minus a veritate aberrabunt, quo ulterius progredientur.

\originalsection{80}Consideremus jam laminam elasticam nusquam fixam, sed liberam vel plano politissimo incumbentem, vel remota gravitate, in spatio vacuo versantem. Facile autem patet hujusmodi laminam motum oscillatorium recipere posse, dum lamina $acb$ sese incurvando alternatim cis \& ultra statum quietis $AB$ excurrit. \marginnote{\hyperlink{addfig:22}{Fig.~22}}Motus igitur iste oscillatorius simili modo, quo in casu præcedente, definiri poterit, dummodo calculus debito modo ad hunc casum accommodetur. Sit igitur $acb$ figura laminæ incurvatæ quam inter oscillandum obtinet, at $ACB$ sit situs ejusdem laminæ in statu æquilibrii, per quem in quavis oscillatione transit. Ponatur, ut ante, longitudo laminæ $AB=a$, ejus elasticitas absoluta $=Ekk$, atque pondus seu massa $=M$. \sourcepage{296}Deinde sit abscissa $AP=x$, applicata $PM=y$; arcus $aM=s$, qui cum abscissa $x$ confundetur; ita ut statui queat $ds=dx$; ex quo radius osculi in $M$ orietur $=\frac{dy^2}{ddy}=R$. Sit autem porro applicata prima $Aa=b$. His positis, ratiocinium ut ante instituendo, ad eandem pervenietur æquationem $\frac{Ekk}R=\frac M{af}\int dx\int ydx=\frac{Ekkddy}{dx^2}$.

\originalsection{81}Si igitur ponamus $\frac{Ekkaf}M=c^4$, ubi $f$ ut ante exprimit longitudinem penduli simplicis isochroni; habebitur, integrando, pro curva hæc æquatio
\[y=Ae^{\frac xc}+Be^{-\frac xc}+C\sin.\frac xc+D\cos.\frac xc;\]
quæ ad præsentem casum ita accommodabitur. Primo, si ponatur $x=0$; fieri debet $y=b$; unde fit
\[b=A+B+D.\]
Secundo, cum sit $\frac{c^4ddy}{dx^2}=\int dx\int ydx$; posito $x=0$, fieri debet $\frac{ddy}{dx^2}=0$, unde prodit
\[0=A+B-D.\]
Tertio, cum sit $\frac{c^4d^3y}{dx^3}=\int ydx$; posito $x=0$, fieri quoque debet $\frac{d^3y}{dx^3}=0$, unde nascitur:
\[0=A-B-D.\]
Quarto, si ponatur $x=a$, evanescere debet $\int ydx$, seu $\frac{d^3y}{dx^3}$; propterea quod $\int ydx$ exprimit summam omnium virium laminam in directione ad axem $AB$ normali trahentium, quæ summa si non esset $=0$, ipsa lamina motu locali promoveretur, contra institutum; erit ergo, ob hanc rationem,
\[0=Ae^{\frac ac}-Be^{-\frac ac}-C\cos.\frac ac+D\sin.\frac ac.\]
\sourcepage{297}Quinto, quia lamina in extremitate $B$ est libera, ibi curvaturam nullam habere poterit; eritque ideo, posito $x=a$, quoque $\frac{ddy}{dx^2}=0$; unde erit
\[0=Ae^{\frac ac}+Be^{-\frac ac}-C\sin.\frac ac-D\cos.\frac ac.\]
His igitur quinque conditionibus in computum ductis, non solum quatuor constantes $A$, $B$, $C$ \& $D$ determinabuntur; sed etiam fractionis $\frac ac$ valor reperietur; ex quo proinde longitudo penduli simplicis isochroni $f$ innotescet.

\originalsection{82}Ex harum æquationum secunda \& tertia, obtinetur $D=A+B$, \& $C=A-B$, qui in sequentibus substituti præbebunt
\[\begin{aligned}
0&=Ae^{\frac ac}-Be^{-\frac ac}-(A-B)\cos.\frac ac+(A+B)\sin.\frac ac,\\
0&=Ae^{\frac ac}+Be^{-\frac ac}-(A-B)\sin.\frac ac-(A+B)\cos.\frac ac;
\end{aligned}\]
ex quibus reperitur:
\[\frac AB=\frac{e^{-\frac ac}-\cos.\frac ac-\sin.\frac ac}{e^{\frac ac}-\cos.\frac ac+\sin.\frac ac}=\frac{-e^{-\frac ac}-\sin.\frac ac+\cos.\frac ac}{e^{\frac ac}-\sin.\frac ac-\cos.\frac ac};\]
ex qua æqualitate elicitur ista æquatio
\[0=2-e^{\frac ac}\cos.\frac ac-e^{-\frac ac}\cos.\frac ac;\quad\text{seu }e^{\frac ac}=\frac{1\pm\sin.\frac ac}{\cos.\frac ac}\]
unde sequentes formabuntur æquationes
\sourcepage{298}
\[\begin{array}{rl}
\mathrm{I}.&\frac ac=\frac12\pi-\varphi=l\operatorname{tang}.\frac12\varphi,\quad\text{quæ dat }\frac ac=0.\\
&\text{pro situ laminæ naturali}\\
\mathrm{II}.&\frac ac=\frac12\pi-\varphi=l\cot.\frac12\varphi\\
\mathrm{III}.&\frac ac=\frac32\pi+\varphi=l\cot.\frac12\varphi\\
\mathrm{IV}.&\frac ac=\frac52\pi-\varphi=l\cot.\frac12\varphi\\
\mathrm{V}.&\frac ac=\frac72\pi+\varphi=l\cot.\frac12\varphi\\
\mathrm{VI}.&\frac ac=\frac92\pi-\varphi=l\cot.\frac12\varphi\\
&\&c.
\end{array}\]

\originalsection{83}Hæ æquationes iterum indicant innumerabiles oscillationum modos, in quorum secundo lamina semel tantum axem $AB$ intersecabit, in tertio bis, in quarto ter, in quinto quater, \& ita porro. Ex quibus intelligitur modos secundum, quartum, sextum \&c. ad præsens institutum non esse accommodatos. Quoniam enim in his numerus intersectionum est impar; laminæ situs inter oscillandum in secundo foret talis, qualem \hyperlink{addfig:23}{Figura 23} repræsentat; in quo quamvis summa virium sollicitantium per totam laminam evanescat, tamen ab iis lamina circa punctum medium $C$ motum rotatorium acquireret: quia vires utrique semissi $aC$ \& $bC$ applicatæ ad eundem laminæ motum rotatorium inducendum conspirarent. Quam ob causam, cum omnino motus rotatorius excludi debeat, figura laminæ, quam inter oscillandum induit, ita debet esse comparata, ut non solum virium sollicitantium toti laminæ applicatarum sit $=0$, sed etiam ut earum summa momentorum evanescat; quod obtinetur si curva in puncto medio $c$, diametro $cC$ sit prædita. \marginnote{\hyperlink{addfig:22}{Fig.~22}}Quod evenit si curva axem $AB$ vel in duobus, vel in quatuor, vel generatim in punctorum numero pari secet; ex quo æquationes tertia, quinta, septima \&c. Solutiones tantum convenientes præbebunt.

\originalsection{84}Hæc ipsa limitatio in ipsa Problematis propositione
\sourcepage{299}contenta reperietur, si ejusmodi tantum curvas admittamus, quæ rectam $Cc$ habeant diametrum, seu in quibus valor ipsius $y$ prodeat idem, si loco $x$ scribatur $a-x$. Ponamus ergo in æquatione generali $a-x$ loco $x$: atque prodibit
\[\begin{aligned}
y&=Ae^{\frac ac}e^{-\frac xc}+Be^{-\frac ac}e^{\frac xc}+C\sin.\frac ac.\cos.\frac xc-C\cos.\frac ac.\sin.\frac xc\\
&\qquad+D\cos.\frac ac.\cos.\frac xc+D\sin.\frac ac.\sin.\frac xc,
\end{aligned}\]
quæ cum congruere debeat cum æquatione
\[y=Ae^{\frac xc}+Be^{-\frac xc}+C\sin.\frac xc+D\cos.\frac xc;\]
fiet $Ae^{\frac ac}=B$, $C(1+\cos.\frac ac)=D\sin.\frac ac$, \& $C\sin.\frac ac=D(1-\cos.\frac ac)$; quarum duæ posteriores congruunt. Cum ergo sit $\frac AB=e^{-\frac ac}$, hoc valore cum superioribus comparato prodibit:
\[e^{-\frac ac}-\cos.\frac ac-\sin.\frac ac=1-e^{-\frac ac}\cos.\frac ac+e^{-\frac ac}\sin.\frac ac,\]
seu
\[e^{-\frac ac}=\frac{1+\cos.\frac ac+\sin.\frac ac}{1+\cos.\frac ac-\sin.\frac ac}=\frac{1+\sin.\frac ac}{\cos.\frac ac}=\frac{\cos.\frac ac}{1-\sin.\frac ac}.\]

\originalsection{85}Erit ergo $e^{\frac ac}=\frac{1-\sin.\frac ac}{\cos.\frac ac}$: sicque in æquatione prius inventa $e^{\frac ac}=\frac{1\pm\sin.\frac ac}{\cos.\frac ac}$; semissis tantum casuum supra exhibitorum, scilicet ii qui sunt numeris imparibus, præsens Problema resolvent. Quare cum prima æquatio contineat \sourcepage{300}laminæ statum naturalem, omnes oscillationum modi in sequentibus æquationibus continebuntur:
\[\begin{array}{rl}
\mathrm{I}.&\frac ac=\frac32\pi+\varphi=l\cot.\frac12\varphi\\
\mathrm{II}.&\frac ac=\frac72\pi+\varphi=l\cot.\frac12\varphi\\
\mathrm{III}.&\frac ac=\frac{11}2\pi+\varphi=l\cot.\frac12\varphi\\
&\&c.
\end{array}\]
Æquationum ergo harum prima præbebit primum eumque principalem oscillandi modum, pro quo valor anguli $\varphi$ simili modo, quo supra, per approximationem reperietur. Limites autem anguli $\varphi$ mox colliguntur esse $1^\circ$, $0'$, $40''$ \& $1^\circ$, $1'$, $0''$, ex quibus per sequentem calculum verus ipsius $\varphi$ valor eruitur.
\[\renewcommand{\arraystretch}{1.12}\begin{array}{r@{\;=\;}r|r}
\varphi&1^\circ,0',40''&1^\circ,1',0''\\
\text{seu}&3640''&3660''\\
\log.&3{,}5611013836&3{,}5634810854\\
\text{subtr.}&5{,}3144251332&5{,}3144251332\\\hline
l\varphi&8{,}2466762504&8{,}2490559522\\
\varphi&0{,}0176472180&0{,}0177441807\\
\frac32\pi&4{,}7123889804&4{,}7123889804\\\hline
\frac ac&4{,}7300361984&4{,}7301331611\\\hline
\frac12\varphi&30',20''&30',30''\\
v&2{,}0543424742&2{,}0519626482\\
lv&0{,}3126728453&0{,}3121694510\\
\text{add.}&0{,}3622156886&0{,}3622156886\\\hline
lu&0{,}6748885339&0{,}6743851396\\
u&4{,}7302983543&4{,}7248186037\\\hline
\text{Error.}&+636341&+53145574\\
\multicolumn{1}{r}{}&&636341\\\cline{3-3}
\text{diff.}&&52509233
\end{array}\]
\sourcepage{301}Hinc intelligitur verum valorem ipsius $\varphi$ non intra istos limites contineri, sed aliquantulum esse minorem quam $1^\circ$, $0'$, $40''$. Nihilo vero minus is ex his erroribus reperietur. Sit enim $\varphi=1^\circ$, $0'$, $40''-n''$; erit $20'':52509233=n'':636341$; unde reperitur $n=\frac{2423}{10000}$, ita ut sit
\[\varphi=1^\circ,0',39\frac{7576''}{10000}.\]
Cum ergo sit $\varphi=3639{,}7576''$ erit
\[\begin{array}{r@{\;=\;}r}
l\varphi&3{,}5610724615\\
\text{subtr.}&5{,}3144251332\\\cline{2-2}
\multicolumn{1}{r}{}&8{,}2466473283\\
\varphi&0{,}0176460428\\
\frac32\pi&4{,}7123889804\\\cline{2-2}
\frac ac&4{,}7300350232.
\end{array}\]

\originalsection{86}Sit hic numerus $=m$, erit, ob $c^4=\frac{Ekkaf}M$, $a^4=\frac{m^4.Ekk.af}M$, \& $f=\frac{a^4}{m^4}.\frac1{Ekk}.\frac Ma$. Unde pari modo numerus oscillationum ab hac lamina uno minuto secundo editarum, erit $=\frac{mm}{aa}\sqrt{g.Ekk.\frac aM}$, existente $g=3{,}16625$ ped. Rhen. Quod si ergo eadem lamina, nunc altero termino $B$ muro infixo, nunc libera ad sonum edendum incitetur, erunt soni inter se ut $nn$ ad $mm$, hoc est ut quadrata numerorum $1{,}8751040813$ \& $4{,}7300350232$, hoc est ut $1$ ad $6{,}363236$. Ratio ergo horum sonorum erit proxime ut $11$ ad $70$: horum ergo sonorum intervallum constituet duas octavas, cum quinta \& hemitonio. Sin autem posterior lamina libera duplo longior capiatur quam prior fixa, intervallum sonorum erit fere sexta minor.

\originalsection{87}Invento hoc valore fractionis $\frac ac$; æquatio pro curva, \sourcepage{302}quam lamina inter oscillandum format, hactenus indeterminata poterit determinari. Cum enim sit $e^{\frac ac}=\frac{1-\sin.\frac ac}{\cos.\frac ac}$, erit $B=\frac{1-\sin.\frac ac}{\cos.\frac ac}A$, \& $C=A-B=A(\cos.\frac ac+\sin.\frac ac-1):\cos.\frac ac$, \& $D=A+B=A(\cos.\frac ac-\sin.\frac ac+1):\cos.\frac ac$. Jam est $b=A+B+D=2D=2A(\cos.\frac ac-\sin.\frac ac+1):\cos.\frac ac$; unde fit
\[\begin{aligned}
A&=\frac{b\cos.\frac ac}{2(\cos.\frac ac-\sin.\frac ac+1)}=\frac{b(1+\sin.\frac ac-\cos.\frac ac)}{4\sin.\frac ac},\\
B&=\frac{b(1-\sin.\frac ac)}{2(\cos.\frac ac-\sin.\frac ac+1)}=\frac{b(-1+\sin.\frac ac+\cos.\frac ac)}{4\sin.\frac ac},\\
C&=\frac{b(-1+\sin.\frac ac+\cos.\frac ac)}{2(\cos.\frac ac-\sin.\frac ac+1)}=\frac{b(1-\cos.\frac ac)}{2\sin.\frac ac},\\
D&=\frac b2=\frac{b\sin.\frac ac}{2\sin.\frac ac}.
\end{aligned}\]
His substitutis oritur hæc æquatio:
\[\begin{aligned}
\frac yb&=\frac{e^{\frac xc}\cos.\frac ac+e^{-\frac xc}(1-\sin.\frac ac)}{2(1-\sin.\frac ac+\cos.\frac ac)}\\
&\qquad+\frac{(1-\cos.\frac ac)\sin.\frac xc+\sin.\frac ac\cos.\frac xc}{2\sin.\frac ac}.
\end{aligned}\]

\originalsection{88}Quia autem recta $Cc$ est curvæ diameter, ponatur abscissa a puncto medio $C$ sumpta $CP=z$, erit $x=\frac12a-z$.
\sourcepage{303}Unde fit $e^{\frac xc}=e^{\frac a{2c}}e^{-\frac zc}=e^{-\frac zc}\sqrt{\frac{1-\sin.\frac ac}{\cos.\frac ac}}$, \& $e^{-\frac xc}=e^{\frac zc}\sqrt{\frac{\cos.\frac ac}{1-\sin.\frac ac}}$; ex quo erit
\[\frac{Ae^{\frac xc}+Be^{-\frac xc}}b=\frac{(e^{\frac zc}+e^{-\frac zc})\sqrt{\cos.\frac ac(1-\sin.\frac ac)}}{2(1-\sin.\frac ac+\cos.\frac ac)}=\frac{e^{\frac zc}+e^{-\frac zc}}{2(e^{\frac a{2c}}+e^{-\frac a{2c}})}.\]
Tum vero erit $(1-\cos.\frac ac)\sin.\frac xc+\sin.\frac ac\cos.\frac xc=\sin.\frac xc+\sin.\frac{a-x}c=\sin.(\frac a{2c}-\frac zc)+\sin.(\frac a{2c}+\frac zc)=2\sin.\frac a{2c}\cos.\frac zc$; quibus substitutis, oritur hæc æquatio:
\[\frac{2y}b=\frac{e^{\frac zc}+e^{-\frac zc}}{e^{\frac a{2c}}+e^{-\frac a{2c}}}+\frac{\cos.\frac zc}{\cos.\frac a{2c}};\]
quæ est forma simplicissima, qua natura curvæ $aMcb$ exprimi potest: manifestum autem est, sive $z$ sumatur affirmative, sive negative, eundem esse proditurum valorem applicatæ $y$. Est vero $e^{\frac a{2c}}+e^{-\frac a{2c}}=\frac{2\cos.\frac a{2c}}{\sqrt{\cos.\frac ac}}$. Invenimus autem angulum $\frac ac=271^\circ$, $0'$, $39\frac34''$.

\originalsection{89}Si jam ponatur $z=0$; præbebit $y$ valorem applicatæ $Cc$; erit ergo $\frac{2.Cc}b=\frac{2\sqrt{\cos.\frac ac}}{2\cos.\frac a{2c}}+\frac1{\cos.\frac a{2c}}$ seu
\sourcepage{304}
\[\frac{Cc}{Aa}=\frac{1+\sqrt{\cos.\frac ac}}{2\cos.\frac a{2c}}=\frac12\sec.\frac a{2c}+\frac12\sec.\frac a{2c}\sqrt{\cos.\frac ac}.\]
At est $\cos.\frac ac=\sin.1^\circ$, $0'$, $39\frac34''$ \& $\cos.\frac a{2c}=\sin.45^\circ$, $30'$, $19\frac78''$. Hinc reperitur $\frac{Cc}{Aa}=0{,}607815$. Deinde si ponatur $y=0$, reperientur puncta $E$ \& $F$ quibus curva axem intersecat. Erit ergo
\[e^{\frac zc}+e^{-\frac zc}=-\frac{\cos.\frac zc}{\cos.\frac a{2c}}(e^{\frac ac}+e^{-\frac a{2c}})=\frac{2\cos.\frac zc}{\sqrt{\cos.\frac ac}},\]
ex qua per approximationes reperitur
\[\frac{CE}{CA}=0{,}551685,\quad\&\quad\frac{AE}{AC}=0{,}448315.\]
Dum ergo lamina oscillationes peragit, hæc puncta $E$ \& $F$ restabunt immobilia; ex quo hujusmodi motus oscillatorius, qui alias vix actu produci posse videatur, facile produci poterit. Si enim lamina in punctis $E$ \& $F$ hoc modo definitis figatur, tum perinde oscillabitur ac si penitus esset libera.

\originalsection{90}Si eodem modo tractetur æquationum supra inventarum secunda $\frac ac=\frac{7\pi}2+\varphi=l\cot.\varphi$; quo quidem casu reperietur proxime $\varphi=0$; tum prodibit secundus modus, quo lamina libera vibrationes absolvere potest, secando scilicet axem $AB$ in quatuor punctis; ideoque lamina perinde oscillabitur, ac si in his quatuor punctis esset fixa. Vicissim ergo, si lamina in his quatuor punctis, vel eorum duobus tantum quibusvis figatur; tum, eodem modo oscillabitur ac si esset libera; sonum autem edet multo auctiorem; quippe qui ad sonum præcedentem modo editum rationem tenebit fere ut $7^2$ ad $3^2$, hoc est, intervallum erit duarum octavarum cum quarta \& hemitonii semisse.
\sourcepage{305}
Tertius oscillandi modus, quo est $\frac ac=\frac{11}2\pi+\varphi=l\cot.\varphi$, habebit sex curvæ $acb$ intersectiones cum axe $AB$; sonusque edetur plus una octava cum tertia minore acutior; huncque sonum lamina edet si in duobus illorum sex punctorum figatur. Hinc patet quam varii soni ab eadem lamina, prout in duobus punctis diversimode figitur, edi queant; \& nisi puncta bina, quibus infigitur, congruant cum intersectionibus in modo primo, vel secundo, vel tertio, atque adeo oscillationes sese ad modorum aliquem sequentium, vel etiam ad infinitesimum componant, tum sonum fore tantopere acutum, ut percipi omnino nequeat, seu quod eodem redit, lamina motum oscillatorium prorsus recipere non poterit: vel saltem, instar cordæ vibrantis, cui ponticulus ita subjicitur ut partes nullam inter se teneant rationem rationalem, sonus minus distinctus producetur.

\originalsection{90bis}\footnote{Original section number 90 is repeated in the source. The suffix bis distinguishes the second occurrence.}Infixa nunc sit lamina elastica in utroque termino $A$ \& $B$; ita tamen ut tangentes curvæ in his punctis non determinentur. Ad hunc scilicet casum in experimentis producendum, laminæ in utroque termino infigantur tenuissimi aculei $A\alpha$, $B\beta$, qui parieti infixi reddant laminæ extremos terminos $A$ \& $B$ immobiles. Ad motum oscillatorium hujus laminæ elasticæ investigandum, ponatur, ut ante, elasticitas absoluta laminæ $=Ekk$, longitudo $AB=a$, \& pondus $=M$, atque longitudo penduli simplicis isochroni $=f$. Sit $AMB$ figura curvilinea, quam lamina inter oscillandum induit, ac ponatur abscissa $AP=AM=x$, applicata $PM=y$, \& radius osculi in $M=R$. Sit porro $P$ vis, quam aculeus $A\alpha$ sustinet in directione $A\alpha$, \& quia vis, qua elementum $Mm$ in directione $M\mu$ urgeri debet quo lamina in hoc statu conservetur, est $=\frac{Mydx}{af}$; erit, per Regulas supra descriptas, æquatio pro curva hæc
\[\frac{Ekk}{R}=Px-\frac M{af}\int dx\int ydx.\]
Est vero $R=-\frac{dx^2}{ddy}$; quia curva versus axem est concava;
\sourcepage{306}
unde fit $\frac{Ekkddy}{dx^2}=\frac M{af}\int dx\int ydx-Px$. Facto ergo $x=0$, erit radius osculi $R$ in $A$ infinitus, ideoque $ddy=0$.\marginnote{\hyperlink{addfig:24}{Fig.~24.}}

\originalsection{91}Si hæc æquatio bis differentietur; prodibit eadem æquatio, quam pro casibus præcedentibus invenimus,
\[Ekkd^4y=\frac M{af}ydx^4.\]
Quod si ergo ponatur $\frac{Ekkaf}{M}=c^4$, erit æquatio integralis
\[y=Ae^{\frac xc}+Be^{-\frac xc}+C\sin.\frac xc+D\cos.\frac xc.\]
Ad quam determinandam, ponatur $x=0$, \& quia simul $y$ evanescere debet, erit $0=A+B+D$. Secundo, ponatur $x=a$, \& quia pariter fieri debet $y=0$, erit
\[0=Ae^{\frac ac}+Be^{-\frac ac}+C\sin.\frac ac+D\cos.\frac ac.\]
Tertio, quia $\frac{ddy}{dx^2}$ evanescere debet, posito \& $x=0$ \& $x=a$; fiet
\[\begin{aligned}0&=A+B-D,\\
0&=Ae^{\frac ac}+Be^{-\frac ac}-C\sin.\frac ac-D\cos.\frac ac.\end{aligned}\]
Jam æquationes $0=A+B-D$ \& $0=A+B+D$ dant $D=0$, \& $B=-A$; qui valores in reliquis duabus æquationibus substituti præbent
\[\begin{aligned}0&=A(e^{\frac ac}-e^{-\frac ac})+C\sin.\frac ac,\\0&=A(e^{\frac ac}-e^{-\frac ac})-C\sin.\frac ac;\end{aligned}\]
quibus satisfieri nequit, nisi sit $A=0$, quia non potest esse $e^{\frac ac}=e^{-\frac ac}$, præter casum $\frac ac=0$; tum vero esse debet $C\sin.\frac ac=0$. Hic cum nequeat poni $C=0$, quia motus oscillatorius foret nullus, erit $\sin.\frac ac=0$, ideoque vel
\sourcepage{307}
$\frac ac=\pi$, vel $\frac ac=2\pi$, \&c. unde iterum infiniti diversi oscillationum oriuntur modi, prout curva $AMB$ axem vel nusquam præter terminos $A$ \& $B$ secat, vel in uno, vel in duobus, vel in pluribus punctis; uti colligitur ex æquatione $y=C\sin.\frac xc$. Puncta intersectionum autem quotcunque fuerint, æqualibus intervallis inter se distabunt.

\originalsection{93}\footnote{The source proceeds from section 91 to section 93; no section numbered 92 is printed.}Cum igitur pro primo ac principali oscillandi modo sit $\frac ac=\pi$, erit $a^4=\pi^4c^4=\pi^4\times Ekk\times\frac aM\times f$, unde fit $f=\frac{a^4}{\pi^4}\times\frac1{Ekk}\times\frac Ma$; Quare ratione longitudinis laminæ, soni iterum tenebunt rationem reciprocam duplicatam longitudinum. Sonus autem hujus laminæ hoc modo editus se habebit ad sonum ejusdem laminæ, si altero termino $B$ muro esset infixa, uti $\pi\pi$ ad quadratum numeri $1{,}8751040813$, hoc est ut $2{,}807041$ ad $1$, seu in numeris minimis ut $57$ ad $160$, quod intervallum est octava cum tritono fere. Si oscillationes se ad secundum modum, quo est $\frac ac=2\pi$, componant; sonus fiet duplici octava acutior; sin sit $\frac ac=3\pi$, sonus acutior fiet tribus octavis cum tono majore, quam casu quo $\frac ac=\pi$; \& ita porro. Quæ quo facilius ad experimenta revocari queant; notandum est oscillationes hic quam-minimas poni, ita ut nulla laminæ elongatione sit opus. Quare, ne tenacitas laminæ, quæ etiam minimæ extensioni, sine qua oscillationes istæ peragi nequeunt, reluctatur, hic alterationem afferat; cuspides illæ ita debent constitui, ut tantilla extensio non impediatur: quod evenit si plano politissimo incumbant. Sic lamina elastica $AB$ in $A$ \& $B$ cuspidibus $A\alpha$ \& $B\beta$ munita, si cuspides speculo imponantur, sonum calculo conformem edet.
\sourcepage{308}
\originalsection{94}Hoc casu expedito, istam de laminis elasticis tractationem claudat motus oscillatorius laminæ elasticæ, utroque termino $A$ \& $B$ muro infixæ, ita ut inter oscillandum puncta $A$ \& $B$ non solum maneant immota, sed etiam recta $AB$ perpetuo sit tangens curvæ $AMB$ in punctis $A$ \& $B$. Hic ergo iterum cavendum est, ut obices terminos $A$ \& $B$ comprehendentes non sint adeo firmi, sed tantillam extensionem quanta ad curvaturam requiritur, permittant. Quæcunque ergo sint vires in terminis $A$ \& $B$ ad laminam continendam requisitæ; ad sequentem pervenietur æquationem differentialem quarti ordinis $Ekkd^4y=\frac M{af}ydx^4$; cujus, si ponatur $\frac{Ekkaf}{M}=c^4$, integralis erit ut supra,
\[y=Ae^{\frac xc}+Be^{-\frac xc}+C\sin.\frac xc+D\cos.\frac xc.\]\marginnote{\hyperlink{addfig:25}{Fig.~25.}}

\originalsection{95}Constantes $A$, $B$, $C$, \& $D$ autem ita debent esse comparatæ, ut posito $x=0$, non solum $y$ evanescat, sed etiam fiat $dy=0$, quia in $A$ curva ab axe $AB$ tangitur. Hoc idem utrumque vero evenire debet, si ponatur $x=a$; unde istæ quatuor æquationes nascentur:
\[\begin{aligned}
\mathrm I.\quad0&=A+B+D,\\
\mathrm {II}.\quad0&=A-B+C,\\
\mathrm {III}.\quad0&=Ae^{\frac ac}+Be^{-\frac ac}+C\sin.\frac ac+D\cos.\frac ac,\\
\mathrm {IV}.\quad0&=Ae^{\frac ac}-Be^{-\frac ac}+C\cos.\frac ac-D\sin.\frac ac.
\end{aligned}\]
Ex harum æquationum prima \& secunda oritur $C=-A+B$, \& $D=-A-B$, qui valores in reliquis duabus substituti dabunt
\[\begin{aligned}
0&=Ae^{\frac ac}+Be^{-\frac ac}-(A-B)\sin.\frac ac-(A+B)\cos.\frac ac,\\
0&=Ae^{\frac ac}-Be^{-\frac ac}-(A-B)\cos.\frac ac+(A+B)\sin.\frac ac.
\end{aligned}\]
\sourcepage{309}
quarum summa ac differentia est,
\[\begin{aligned}
0&=Ae^{\frac ac}+B\sin.\frac ac-A\cos.\frac ac,\\
\text{seu}\quad\frac AB&=\frac{\sin.\frac ac}{\cos.\frac ac-e^{\frac ac}};\\
0&=Be^{-\frac ac}-A\sin.\frac ac-B\cos.\frac ac,\\
\text{seu}\quad\frac AB&=\frac{e^{-\frac ac}-\cos.\frac ac}{\sin.\frac ac}.
\end{aligned}\]
unde fit $2=(e^{\frac ac}+e^{-\frac ac})\cos.\frac ac$, seu
\[e^{\frac ac}=\frac{1\pm\sin.\frac ac}{\cos.\frac ac}.\]
Quæ æquatio, quia congruit cum ea quam §.~81 invenimus, sequentes Solutiones numero infinitæ satisfacient:
\[\begin{aligned}
\mathrm I.\quad\frac ac&=\tfrac12\pi-\varphi=l\cot.\tfrac12\varphi,\\
\mathrm {II}.\quad\frac ac&=\tfrac32\pi+\varphi=l\cot.\tfrac12\varphi,\\
\mathrm {III}.\quad\frac ac&=\tfrac52\pi-\varphi=l\cot.\tfrac12\varphi.\\
&\text{\&c.}
\end{aligned}\]
\originalsection{96}Harum æquationum primæ satisfieri nequit, nisi sit $\varphi=90^\circ$, ideoque $\frac ac=0$; unde primus oscillandi modus oritur ex æquatione $\frac ac=\frac32\pi+\varphi=l\cot.\frac12\varphi$; quæ cum jam supra sit tractata, erit $\frac ac=4{,}7300350232$. Quamobrem lamina elastica, cujus uterque terminus parieti infixus tenetur, perinde vibrationes suas peraget, ac si esset omnino libera. Hæc
\sourcepage{310}
autem convenientia tantum ad primum oscillandi modum spectat; secundus enim oscillandi modus, quo est $\frac ac=\frac52\pi-\varphi=\log.\cot.\frac12\varphi$, atque lamina axem $AB$ inter oscillandum in uno puncto intersecat, in lamina libera sui parem non habet; tertius autem modus laminæ utrinque infixæ congruet cum modo secundo laminæ liberæ, atque ita porro.

\originalsection{97}Hæc duo postrema oscillationum genera, ob causam allatam, non congrue per experimenta explorari possunt: primum autem non solum ad experimenta instituenda maxime est aptum; sed etiam adhiberi potest ad elasticitatem absolutam cujusque laminæ propositæ, quam per $Ekk$ indicavimus, investigandam. Quod si enim sonus notetur, quem hujusmodi lamina altero termino muro infixa edit, eique in corda consonus efficiatur, simul numerus oscillationum uno minuto secundarum editarum cognoscetur. Qui si æqualis ponatur expressioni $\frac{nn}{aa}\sqrt{g\cdot Ekk\cdot\frac aM}$, ob numerum $n$ cognitum, \& quantitates $g$, $a$, \& $M$ per dimensiones inventas, reperietur valor expressionis $Ekk$; sicque elasticitas absoluta innotescet; quæ cum ea quam supra ex incurvatione reperire docuimus, comparari potest.
